\documentclass[11pt]{article}

\usepackage[T1]{fontenc}
\usepackage{lmodern}
\usepackage{microtype}
\usepackage[a4paper,margin=28mm]{geometry}
\emergencystretch=2.5em
\usepackage{amsmath,amssymb,amsthm,mathtools,mathrsfs}
\usepackage{enumitem}
\usepackage{booktabs}
\usepackage{longtable}
\usepackage{array}
\usepackage{xcolor}
\usepackage{xurl}
\usepackage{aliascnt}
\usepackage{hyperref}
\usepackage[nameinlink,capitalise,noabbrev]{cleveref}

\hypersetup{
  colorlinks=true,
  linkcolor=blue!55!black,
  citecolor=blue!55!black,
  urlcolor=blue!55!black,
  pdftitle={Hybrid arithmetic holonomy and twenty-two individual p-adic zeta values},
  pdfauthor={Christopher D. Long}
}

\newtheorem{theorem}{Theorem}[section]
\newaliascnt{proposition}{theorem}
\newtheorem{proposition}[proposition]{Proposition}
\aliascntresetthe{proposition}
\newaliascnt{lemma}{theorem}
\newtheorem{lemma}[lemma]{Lemma}
\aliascntresetthe{lemma}
\newaliascnt{corollary}{theorem}
\newtheorem{corollary}[corollary]{Corollary}
\aliascntresetthe{corollary}
\newaliascnt{claim}{theorem}
\newtheorem{claim}[claim]{Claim}
\aliascntresetthe{claim}
\theoremstyle{definition}
\newaliascnt{definition}{theorem}
\newtheorem{definition}[definition]{Definition}
\aliascntresetthe{definition}
\newaliascnt{hypothesis}{theorem}
\newtheorem{hypothesis}[hypothesis]{Hypothesis}
\aliascntresetthe{hypothesis}
\theoremstyle{remark}
\newaliascnt{remark}{theorem}
\newtheorem{remark}[remark]{Remark}
\aliascntresetthe{remark}
\newaliascnt{warning}{theorem}
\newtheorem{warning}[warning]{Warning}
\aliascntresetthe{warning}

\crefname{theorem}{Theorem}{Theorems}
\crefname{proposition}{Proposition}{Propositions}
\crefname{lemma}{Lemma}{Lemmas}
\crefname{corollary}{Corollary}{Corollaries}
\crefname{claim}{Claim}{Claims}
\crefname{definition}{Definition}{Definitions}
\crefname{hypothesis}{Hypothesis}{Hypotheses}
\crefname{remark}{Remark}{Remarks}
\crefname{warning}{Warning}{Warnings}

\newcommand{\Q}{\mathbf Q}
\newcommand{\Z}{\mathbf Z}
\newcommand{\C}{\mathbf C}
\newcommand{\R}{\mathbf R}
\newcommand{\F}{\mathbf F}
\newcommand{\Qp}{\mathbf Q_p}
\newcommand{\dd}{\,\mathrm d}
\newcommand{\ord}{\operatorname{ord}}
\newcommand{\rank}{\operatorname{rank}}
\newcommand{\lcmop}{\operatorname{lcm}}
\newcommand{\Ann}{\operatorname{Ann}}
\newcommand{\gr}{\operatorname{gr}}
\newcommand{\Sym}{\operatorname{Sym}}
\newcommand{\Fil}{\operatorname{Fil}}
\newcommand{\Span}{\operatorname{Span}}
\newcommand{\calE}{\mathcal E}
\newcommand{\calF}{\mathcal F}
\newcommand{\calH}{\mathcal H}
\newcommand{\calK}{\mathcal K}
\newcommand{\calA}{\mathcal A}
\newcommand{\calM}{\mathcal M}
\newcommand{\calL}{\mathcal L}
\newcommand{\calO}{\mathcal O}
\newcommand{\calV}{\mathcal V}
\newcommand{\widehatdeg}{\widehat{\deg}}
\newcommand{\Energy}{\mathscr H}
\newcommand{\LambdaAH}{\Lambda}
\newcommand{\floor}[1]{\left\lfloor #1\right\rfloor}
\newcommand{\ceil}[1]{\left\lceil #1\right\rceil}
\newcommand{\abs}[1]{\left|#1\right|}
\newcommand{\set}[1]{\left\{#1\right\}}
\newcommand{\positivepart}[1]{\left(#1\right)_{+}}
\newcommand{\T}{\mathbb T}

\title{Hybrid arithmetic holonomy and\\
twenty-two individual $p$-adic zeta values}
\author{Christopher D. Long}
\date{Re-audited product-digit repaired research draft --- 24 August 2026}

\begin{document}
\maketitle

\begin{abstract}
We develop a hybrid finite-place refinement of arithmetic holonomy for the
Eisenstein--Eichler connection on the genus-zero modular curves
$X_0(p)$, $p\in\{2,3,5,7\}$.  The denominator staircase of the Eichler
primitive cannot be realized by globally defined scalar solution coordinates:
the nonzero modular weights create a genuine torus obstruction.  We instead
use two different local mechanisms.  At auxiliary primes below a cutoff, the
Hasse invariant algebraizes the unipotent comparison coordinate modulo the
prime and a genuine-source rank--nullity argument produces one
determinant-level saving on a large source kernel.  Above the cutoff, we keep the calculation on the algebraic de Rham
frame torsor; exact divided-Frobenius relations for the raw Eichler rows and a
coefficient-ring Taylor no-backflow lemma recover the Calegari--Dimitrov--Tang
prime window without scalarizing the packet.  The resulting arithmetic cost is
\[
 \tau_{p,s}(\xi)=\frac{2}{(s+1)^2}
 \left(S_{p,s}(\xi)+s I^{s+1}_{\xi}(\xi)\right),
\]
where $S_{p,s}$ is the small-prime Hasse-codimension cost and
$I^{s+1}_{\xi}$ is the published prime-window function.  Combining this with
Calegari's overconvergent Eisenstein family, Buzzard's analytic continuation,
an exact modular Jensen formula, and Bost's slopes inequality proves the
irrationality of at least twenty-two individual Kubota--Leopoldt values:
\[
 \zeta_2(s)\quad(3\le s\le29,\ s\text{ odd}),\qquad
 \zeta_3(s)\quad(3\le s\le11,\ s\text{ odd}),
\]
\[
 \zeta_5(3),\quad \zeta_5(5),\quad \zeta_7(3).
\]
Every numerical inequality is supplied with an independent rational
fixed-point interval certificate.  No opposite Hodge filtration, rational Tate
scalarization, or rowwise transport of literal $n^e$ denominators is used.
\end{abstract}

\medskip
\noindent\textbf{Audit status.}
The proof isolates five published interfaces: Calegari's overconvergent
Eisenstein family, Buzzard's finite-slope continuation theorem, Katz's Hasse
invariant formalism, the de Rham frame-torsor/Gauss--Manin construction of
Graham--Pilloni--Rodrigues Jacinto, and the Bost--Charles/CDT slopes and
prime-window estimates.  The canonical BGG lift, single-action deepest-row
projection, genuine-source Hasse kernel, hybrid local lattices, full-product
pole-grade Cartier argument, collision calculation, and all numerical
applications are proved in the manuscript.  Because the conclusions include
individual irrationality statements beyond the classical cases, independent
specialist review remains appropriate before publication.

\tableofcontents

\section{Introduction and main results}\label{sec:intro}

For an odd integer $s\ge3$, let $\zeta_p(s)$ denote the Kubota--Leopoldt
$p$-adic zeta value in the normalization used by Calegari.  Calegari proved the
irrationality of $\zeta_2(3)$ and $\zeta_3(3)$ in \cite{Calegari2005}; the
irrationality of $\zeta_2(5)$ was later obtained through arithmetic holonomy
and independently by an Ap\'ery-type construction.  General results also give
many irrational values or finite alternatives, but usually do not identify a
large fixed list of individual values; see, for example,
\cite{Lai2025,LaiLupuSprang2025,LaiSprang2025}.  We do not attempt a complete
priority analysis here.  Our theorem is the following fixed list.

\begin{theorem}[Twenty-two individual values]\label{thm:main}
The following $p$-adic numbers are irrational:
\[
\begin{aligned}
 &\zeta_2(s) &&(s=3,5,\ldots,29),\\
 &\zeta_3(s) &&(s=3,5,\ldots,11),\\
 &\zeta_5(3),\ \zeta_5(5),\qquad
 &\zeta_7(3).&
\end{aligned}
\]
\end{theorem}

The proof begins with the negative-weight Eisenstein series
\[
 K_{p,s}(q)=\frac{\zeta_p(s)}2+
 \sum_{n\ge1}\sigma_{-s}^{(p)}(n)q^n,
 \qquad
 \sigma_{-s}^{(p)}(n)=
 \sum_{\substack{d\mid n\\p\nmid d}}d^{-s}.
\]
Under the contrary assumption $\zeta_p(s)\in\Q$, this is an arithmetic
horizontal leaf of a rank-$(s+1)$ Eisenstein--Eichler connection.  Its raw
Euler derivatives have exact integration depths
\[
 0,s,s-1,\ldots,1,
\]
but the positive-depth rows have different modular weights.  The diagonal
torus part of a Hodge-frame change therefore cannot be represented by a
bounded rational matrix in the Hauptmodul.  A scalar filtered realization is
not merely unavailable; it is incompatible with the modular transformation
law.

The replacement is hybrid.  At small auxiliary primes, the Hasse invariant
makes the normalized unipotent comparison coordinate rational modulo the
prime.  A rank--nullity argument on the genuine global source then produces a
kernel of the required codimension, allowing one reciprocal prime to be saved
on
\[
 (s-1)\left(D-\frac{p+1}{12}\ell\right)+O_{p,s}(1)
\]
source directions.  At large auxiliary primes, the scalar pairing is computed
upstairs on the de Rham frame torsor.  The fibre coordinates remain
coefficients, while only the raw scalar Eichler rows undergo the exact Dwork
substitution.  Residue classes modulo $\ell$ prevent Taylor backflow over any
coefficient algebra and recover a degree-$D+O(1)$ prime window.

For $m=s+1$ define
\[
 K_m(\xi)=\floor{\frac{m-1}{\xi}}
\]
and
\begin{equation}\label{eq:I-def-intro}
 I^m_\xi(\xi)=
 \left(m-\frac12\right)H_{K_m(\xi)}-K_m(\xi)\xi
 +\frac{\positivepart{m-(K_m(\xi)+1)\xi}^2}{2(K_m(\xi)+1)}.
\end{equation}
Put
\[
 c_p=\frac{p+1}{12},\qquad
 J_p(\xi)=\int_0^\xi(1-c_px)_+\,\dd x
\]
and
\begin{equation}\label{eq:S-def-intro}
 S_{p,s}(\xi)=s^2\xi-(s-1)J_p(\xi).
\end{equation}
Thus, when $c_p\xi<1$,
\[
 S_{p,s}(\xi)=s^2\xi-(s-1)
 \left(\xi-\frac{p+1}{24}\xi^2\right).
\]
The main arithmetic statement is as follows.

\begin{theorem}[Hybrid arithmetic holonomy]\label{thm:hybrid-holonomy}
Fix $p\in\{2,3,5,7\}$ and an odd $s\ge3$.  Under the assumption
$\zeta_p(s)\in\Q$, let $\boldsymbol\phi$ be admissible adelic continuation
templates for the rank-$m=s+1$ Eisenstein--Eichler leaf, with derivative width
$\LambdaAH(\boldsymbol\phi)$ and Bost--Charles energy
$\Energy(\boldsymbol\phi)$.  For every $1<\xi<m$ one has
\[
 m\left(\LambdaAH(\boldsymbol\phi)-\tau_{p,s}(\xi)\right)
 \le \Energy(\boldsymbol\phi),
\]
where
\begin{equation}\label{eq:tau-main}
 \boxed{
 \tau_{p,s}(\xi)=\frac{2}{(s+1)^2}
 \left(S_{p,s}(\xi)+sI^{s+1}_\xi(\xi)\right).}
\end{equation}
\end{theorem}

For the modular templates used below,
\[
 \Lambda_p(Y)=\frac{12\log p}{p-1}-2\pi Y
\]
and
\[
 \Energy_p(Y)=\Lambda_p(Y)+C_p(Y),
\]
where
\begin{equation}\label{eq:C-intro}
 C_p(Y)=4\pi
 \sum_{\substack{c\ge1\\p\mid c\\cY<1}}
 \frac{\varphi(c)}{c^2}
 \left(\arccos(cY)-cY\sqrt{1-c^2Y^2}\right).
\end{equation}
Consequently rationality would force
\begin{equation}\label{eq:margin-intro}
 \mathfrak M_{p,s}(\xi,Y):=
 s\Lambda_p(Y)-(s+1)\tau_{p,s}(\xi)-C_p(Y)\le0.
\end{equation}
For the twenty-two pairs in \cref{thm:main}, exact rational witnesses
$\xi,Y$ make the left side strictly positive.  The smallest independently
certified lower bound is
\[
 \mathfrak M_{5,5}\left(\frac{29}{27},\frac1{16}\right)
 >0.1317993568270.
\]

\begin{warning}[What is not proved or used]\label{warn:no-scalarization}
The argument does not construct globally defined scalar functions whose
ordinary-LCM depths are $0,1,\ldots,s$.  It does not reverse the Hodge flag,
transport literal $n^e$ denominators through a change of local coordinate, or
replace the nonzero-weight Tate trivializations by rational functions of the
Hauptmodul.  The small-prime saving is only one power of the auxiliary prime;
the large-prime theorem is proved directly on the de Rham frame torsor.
\end{warning}

The paper is organized as follows.  \Cref{sec:inputs} states the published
interfaces and fixes notation.  \Cref{sec:packet} constructs the algebraic
Eisenstein--Eichler packet and proves its raw arithmetic, BGG realization,
functional rank, and zero estimate.  \Cref{sec:sources} constructs the global
source and proves the fixed-data descent and single-action projection bridges.  The small-prime Hasse
kernel and large-prime torsor Cartier theorem are developed in
\cref{sec:hybrid-local}.  \Cref{sec:bost} assembles these local estimates by
Bost's slopes inequality.  Analytic continuation and the exact collision
energy are proved in \cref{sec:continuation,sec:collision}; the twenty-two
applications and the interval certificate appear in
\cref{sec:applications,app:certificate}.

\section{Published inputs and notation}\label{sec:inputs}

We use the following published results.

\begin{enumerate}[label=\textup{(E\arabic*)},leftmargin=3em]
\item \textbf{Overconvergent Eisenstein families.}
Write $s=2n+1$ and set $\kappa(a)=a^{-2n}=a^{1-s}$.  In Calegari's notation,
\[
 \sigma_\kappa^*(m)
 =\sum_{\substack{d\mid m\\p\nmid d}}\kappa(d)d^{-1}
 =\sum_{\substack{d\mid m\\p\nmid d}}d^{-s}.
\]
Thus Theorem~2.3 of \cite{Calegari2005} gives an overconvergent eigenform
with the required nonconstant $q$-expansion, and Lemma~2.4 identifies its
constant term as $\zeta_p(s)/2$.  This is the form denoted $K_{p,s}$ below.

\item \textbf{Finite-slope continuation.}
Buzzard's Theorem~5.2 in \cite{Buzzard2003} extends an overconvergent modular
form of integer weight and nonzero $U_p$-eigenvalue across the wide-open
$W_0(p)$.

\item \textbf{Hasse invariants and $q$-expansions.}
We use Katz's geometric Hasse invariant and the $q$-expansion principle from
\cite{Katz1973}.  For every auxiliary prime $\ell\ge5$, the classical
congruences
\[
 E_{\ell-1}\equiv1\pmod\ell,
 \qquad E_{\ell+1}\equiv E_2\pmod\ell
\]
are used in $q$-expansions.

\item \textbf{The de Rham frame torsor.}
Propositions~A.14--A.15 of Graham--Pilloni--Rodrigues Jacinto
\cite{GPRJ2025} provide, at neat level, the algebraic Hodge-compatible
$B$-torsor, its torus-weight decomposition, the filtration induced by the
unipotent radical, and the Gauss--Manin connection.  The descent from one
fixed neat cover to the modular stack and coarse curve, and the uniform
spreading out over \(\Z[\Sigma^{-1}]\), are not quoted from that source; they
are proved internally in \cref{lem:neat-descent}.

\item \textbf{Arithmetic holonomy.}
We use Bost's slopes inequality \cite{Bost2001}; the arithmetic
Hilbert--Samuel source-degree asymptotic together with the Bost--Charles
overflow identities
\cite[Theorem~5.4.1 and Proposition~5.4.2]{BostCharles2022}, in the adelic
assembly of \cite[Theorem~2.6 and Section~2.3]{CDT2025}; and the
Chudnovsky--Osgood estimate and scalar source, pivot-height, and prime-window
formulas of \cite[Theorem~3.2.13 and Section~7.3, especially
(7.3.7) and (7.3.14)--(7.3.16)]{CDT2024}.  The scalar local-denominator
hypothesis is replaced by the hybrid finite-place theorem proved here; once
the interval window is obtained, the CDT prime-number-theorem summation is
used without alteration.
\end{enumerate}

All structural algebraic objects independent of the hypothetical rational
constant $\eta$ are spread over $\Z[\Sigma^{-1}]$ after enlarging a fixed
finite set $\Sigma$ depending on $p$ and $s$.  Under the contradiction
hypothesis $\eta\in\Q$, the finitely many primes dividing its denominator
are treated separately; they contribute only $o(D^2)$.  Such fixed changes
alter neither a $D^2$ arithmetic-degree coefficient nor a leading pivot-height
coefficient.

Let
\[
 f_p(\tau)=
 \left(\frac{\Delta(p\tau)}{\Delta(\tau)}\right)^{1/(p-1)}.
\]
For $p\in\{2,3,5,7\}$ this is a Hauptmodul at the cusp $P=\infty$ and
\begin{equation}\label{eq:f-product}
 f_p(q)=q\prod_{n\ge1}
 \left(1+q^n+\cdots+q^{(p-1)n}\right)^{24/(p-1)}
 \in q+q^2\Z[[q]].
\end{equation}
Its inverse satisfies
\begin{equation}\label{eq:q-inverse}
 q(f_p)\in f_p+f_p^2\Z[[f_p]].
\end{equation}
We usually write $f=f_p$ and
\[
 D=q\frac{\dd}{\dd q}.
\]

For an admissible adelic family of continuation templates
$\boldsymbol\phi=(\phi_v)_v$, each centered at $P$, put
\[
 \LambdaAH(\boldsymbol\phi)=
 \sum_v\log\abs{\phi_v'(0)}_v
\]
and
\[
 \Energy(\boldsymbol\phi)=
 \sum_{v\mid\infty}
 \iint_{\T_v^2}\log\abs{\phi_v(z)-\phi_v(w)}_v\,\dd m(z)\dd m(w)
 +\sum_{v\nmid\infty}
 \sup_{z\in\T_v}\log\abs{\phi_v(z)}_v.
\]
These are normalized as in \cite[Section~7.3]{CDT2024} and
\cite[Section~2]{CDT2025}.

\section{The Eisenstein--Eichler packet}\label{sec:packet}

\subsection{The negative-weight Eisenstein form}

For odd $s\ge3$ define
\begin{equation}\label{eq:J-def}
 J_{p,s}(q)=\sum_{n\ge1}\sigma_{-s}^{(p)}(n)q^n,
 \qquad
 \sigma_{-s}^{(p)}(n)=
 \sum_{\substack{d\mid n\\p\nmid d}}d^{-s},
\end{equation}
and, for a scalar $\eta$,
\[
 K_{p,s,\eta}=J_{p,s}+\eta.
\]
At
\[
 \eta=\zeta_p(s)/2,
\]
input (E1) identifies $K_{p,s,\eta}$ as an overconvergent eigenform of
weight $1-s$.

\begin{lemma}[$U_p$-eigenvalue]\label{lem:Up}
For every $n\ge1$,
\[
 \sigma_{-s}^{(p)}(pn)=\sigma_{-s}^{(p)}(n).
\]
Consequently
\[
 U_pK_{p,s,\zeta_p(s)/2}=K_{p,s,\zeta_p(s)/2}.
\]
\end{lemma}

\begin{proof}
The divisors of $pn$ that are prime to $p$ are precisely the divisors of $n$
that are prime to $p$.  The usual $q$-expansion definition of $U_p$ now gives
the assertion, including the constant term.
\end{proof}

\subsection{Raw integration depths}

For $1\le e\le s$ put
\begin{equation}\label{eq:R-def}
 R_e=D^{s-e}J_{p,s}.
\end{equation}

\begin{lemma}[Raw integration depth]\label{lem:raw-depth}
For every $n\ge1$,
\begin{equation}\label{eq:raw-depth}
 n^e[q^n]R_e
 =\sum_{\substack{d\mid n\\p\nmid d}}(n/d)^s\in\Z.
\end{equation}
Consequently, if $L_N=\lcmop(1,\ldots,N)$, then
\[
 L_N^eR_e(q)\pmod{q^{N+1}}
\]
has integral coefficients.  The same is true after the integral substitution
$q=q(f)$:
\[
 L_N^eR_e(q(f))\pmod{f^{N+1}}\in\Z[[f]]/(f^{N+1}).
\]
\end{lemma}

\begin{proof}
The first identity follows directly from
$[q^n]R_e=n^{s-e}\sigma_{-s}^{(p)}(n)$.  Since $n\mid L_N$ for
$n\le N$, the $q$-integrality follows.  By \eqref{eq:q-inverse}, the
coefficient of $f^k$ after substitution uses only $q^n$ with $n\le k$, and
all composition coefficients are integral.  This proves the final assertion.
\end{proof}

\begin{remark}
The lemma transports a cumulative LCM cutoff.  It does not assert the
coordinatewise identity $n^e[f^n]R_e(f)\in\Z$, which need not be invariant
under an integral tangent-to-identity change of parameter.
\end{remark}

\subsection{Exact divided Frobenius}

\begin{lemma}[Exact Dwork relation]\label{lem:Dwork}
Let $\ell\ne p$ be a rational prime.  For $1\le e\le s$,
\begin{equation}\label{eq:Dwork-q}
 \boxed{R_e(q)-\ell^{-e}R_e(q^\ell)\in\Z_\ell[[q]].}
\end{equation}
After $q=q(f)$,
\begin{equation}\label{eq:Dwork-f}
 R_e(q(f))-\ell^{-e}R_e(q(\varphi_\ell(f)))
 \in\Z_\ell[[f]],
\end{equation}
where
\[
 \varphi_\ell(f)=f(q(f)^\ell).
\]
\end{lemma}

\begin{proof}
If $\ell\nmid n$, every divisor of $n$ is an $\ell$-adic unit.  If
$\ell\mid n$, the $q^n$-coefficient of the left side of
\eqref{eq:Dwork-q} is
\[
 n^{s-e}\left(
 \sigma_{-s}^{(p)}(n)-\ell^{-s}\sigma_{-s}^{(p)}(n/\ell)
 \right).
\]
The term in parentheses is
\[
 \sum_{\substack{d\mid n\\p\nmid d\\\ell\nmid d}}d^{-s},
\]
which is $\ell$-integral.  Substitution of the integral inverse coordinate
gives \eqref{eq:Dwork-f}.
\end{proof}

\begin{lemma}[Coordinate Frobenius]\label{lem:coordinate-Frob}
For every prime $\ell$ outside a fixed exceptional set,
\begin{equation}\label{eq:Frob-expansion}
 \boxed{
 \varphi_\ell(f)=f^\ell+\ell\Delta_\ell(f),
 \qquad
 \Delta_\ell(f)\in f^{\ell+1}\Z_\ell[[f]].}
\end{equation}
\end{lemma}

\begin{proof}
The product \eqref{eq:f-product} is integral and begins with $q$.  Hence
in characteristic $\ell$,
\[
 f(q)^\ell=f(q^\ell).
\]
Substituting $q=q(f)$ gives
$\varphi_\ell(f)\equiv f^\ell\pmod\ell$.  Both sides have order
$\ell$ and the same leading coefficient, so the quotient of their difference
by $\ell$ has order at least $\ell+1$.
\end{proof}

\subsection{The split BGG extension}

Let $\calH$ be the rank-two de Rham bundle of the universal elliptic curve,
with its logarithmic Gauss--Manin connection, and put
\[
 V=\Sym^{s-1}\calH.
\]
Let $L_{\mathrm{top}}\subset V$ be the top oper line and
$L_{\mathrm{bot}}$ the bottom oper quotient.  Locally choose an oper
derivation $\partial$ and an oper frame $v_0,\ldots,v_{s-1}$ such that
\[
 L_{\mathrm{top}}=\langle v_0\rangle,
 \qquad
 L_{\mathrm{bot}}=V/\langle v_0,\ldots,v_{s-2}\rangle,
\]
and
\[
 \nabla_\partial v_j=(s-1-j)v_{j+1},
 \qquad v_s=0.
\]
For a local section $F$ of $L_{\mathrm{bot}}\simeq\omega^{1-s}$,
represented as a scalar in this oper gauge, define
\begin{equation}\label{eq:spl-def}
 \operatorname{spl}_s(F)=
 \sum_{j=0}^{s-1}
 (-1)^j\frac{(s-1)!}{(s-1-j)!}
 \partial^{s-1-j}F\,v_j.
\end{equation}

\begin{lemma}[Canonical BGG lift and normalization]\label{lem:BGG}
The expression \eqref{eq:spl-def} is the unique local section $G$ of $V$
with the following two intrinsic properties:
\begin{enumerate}[label=\textup{(\roman*)}]
\item the image of $G$ in $L_{\mathrm{bot}}$ is
$(-1)^{s-1}(s-1)!F=(s-1)!F$;
\item $\nabla_VG$ takes values in
$L_{\mathrm{top}}\otimes\Omega^1(\log D)$.
\end{enumerate}
Consequently the local expressions glue in the Zariski topology to a
canonical algebraic differential operator
\[
 \operatorname{spl}_s:\omega^{1-s}\longrightarrow V
\]
of order $s-1$.  This is the $\mathrm{GL}_2$ dual-BGG lift in the present
normalization; compare \cite[Proposition~2.4.3]{RodriguezCamargo2023}.  Moreover,
\begin{equation}\label{eq:BGG}
 \nabla_V\operatorname{spl}_s(F)
 =\iota_s(\operatorname{Bol}_sF).
\end{equation}
Here $\iota_s$ is normalized by the oper frame so that its image is the
$v_0$-line, and $\operatorname{Bol}_sF=\partial^sF$ in an oper coordinate.
On the Tate disc this is $D^sF$, with $D=q\,\mathrm d/\mathrm dq$.
\end{lemma}

\begin{proof}
Put $c_j=(-1)^j(s-1)!/(s-1-j)!$.  The image of
\eqref{eq:spl-def} in the bottom quotient is $c_{s-1}F=(s-1)!F$, because
$s-1$ is even.  For $j\ge1$, the coefficient of
$\partial^{s-j}Fv_j$ in $\nabla_\partial\operatorname{spl}_s(F)$ is
\[
 c_j+(s-j)c_{j-1}=0.
\]
Only the $v_0$-term $\partial^sFv_0$ remains, proving existence and
\eqref{eq:BGG} with the displayed normalization.

For uniqueness, let $W=\sum_{j=0}^{s-1}a_jv_j$ have zero image in
$L_{\mathrm{bot}}$ and suppose $\nabla_VW$ lies in the top oper line.  Then
$a_{s-1}=0$, and the vanishing of the $v_j$-coefficient of $\nabla_\partial W$
for $1\le j\le s-1$ gives
\[
 \partial a_j+(s-j)a_{j-1}=0.
\]
Starting with $j=s-1$ and descending gives
$a_{s-2}=\cdots=a_0=0$.  Thus the two intrinsic properties determine the
lift uniquely.  The oper filtration, logarithmic connection, and recursion
above are algebraic, so local formulas on overlapping algebraic oper charts
have those same properties, agree, and glue to an algebraic differential
operator over the field of definition.

For the exact Tate normalization, write
\[
 \delta_k^{[a]}=
 \delta_{k+2(a-1)}\cdots\delta_{k+2}\delta_k.
\]
Urban's iterated Gauss--Manin formula
\cite[Proposition~2.4.1 and Remark~2.4.2]{Urban2013} gives, in the
polynomial realization on the de Rham torsor,
\begin{equation}\label{eq:Urban-iterate}
 \delta_k^{[a]}F
 =\sum_{j=0}^{a}\binom aj
   \left(\prod_{b=1}^{j}(k+a-b)\right)
   X^jD^{a-j}F.
\end{equation}
Although Urban writes the formula in the classical nonnegative-weight range,
its proof is the formal induction from $\delta_k=D+kX$ and
$[D,X]=-X^2$; both sides are polynomial in $k$.  Hence
\eqref{eq:Urban-iterate} remains valid at the negative integer
$k=1-s$.  Taking $a=s-1$, the product in the $j$th coefficient is
$(-1)^j j!$, so \eqref{eq:Urban-iterate} gives exactly the coefficients in
\eqref{eq:spl-def}, with no unrecorded scalar or sign.  Taking instead
$a=s$, every term with $j\ge1$ contains the factor $k+s-1=0$, and hence
\[
 \delta_{1-s}^{[s]}F=D^sF.
\]
The global derivation and weight decomposition on the Hodge-compatible frame
torsor are those of \cite[Propositions~A.14--A.15]{GPRJ2025}; thus this Tate
formula is the torsor realization of the canonical lift just constructed.
\end{proof}

Put
\begin{equation}\label{eq:Eevil}
 E_{p,s+1}^{\mathrm{evil}}:=D^sK_{p,s,\eta}=D^sJ_{p,s}.
\end{equation}
Its $q$-expansion is
\[
 \sum_{n\ge1}
 \left(\sigma_s(n)-\mathbf1_{p\mid n}\sigma_s(n/p)\right)q^n,
\]
so it is the algebraic $p$-stabilized Eisenstein form of weight $s+1$ and is
independent of $\eta$.

\begin{proposition}[Split algebraic Eisenstein--Eichler connection]
\label{prop:split-extension}
On the split vector bundle
\[
 \calE_{p,s}=\calO\oplus V
\]
define
\begin{equation}\label{eq:extension-connection}
 \nabla_{\calE}(a,v)=
 \left(\dd a,\nabla_Vv-a\,\iota_s(E_{p,s+1}^{\mathrm{evil}})\right).
\end{equation}
This is an algebraic regular-singular connection, and on the Tate formal disc
\[
 S_{p,s,\eta}=
 \left(1,\operatorname{spl}_s(K_{p,s,\eta})\right)
\]
is horizontal.
\end{proposition}

\begin{proof}
The connection is algebraic because $E_{p,s+1}^{\mathrm{evil}}$ is algebraic;
integrability is automatic on a curve and regular singularity follows from the
logarithmic Gauss--Manin connection.  By \cref{lem:BGG},
\[
 \nabla_V\operatorname{spl}_s(K_{p,s,\eta})
 =\iota_s(E_{p,s+1}^{\mathrm{evil}}),
\]
which is exactly the horizontality equation for the displayed vector.
\end{proof}

\subsection{Rees--Bol jets}

On the Tate de Rham torsor use Urban's polynomial coordinate $X$ and
operators, with the polynomial $q$-expansion normalization of
\cite[Proposition~2.4.1 and Remark~2.4.2]{Urban2013} and the torsor
formalism of \cite{GPRJ2025},
\[
 \delta_k=D+kX,
 \qquad
 D=q\frac{\partial}{\partial q}-X^2\frac{\partial}{\partial X},
 \qquad[D,X]=-X^2.
\]
Define
\[
 \calK_0=K_{p,s,\eta},
 \qquad
 \calK_j=
 \delta_{1-s+2(j-1)}\cdots\delta_{1-s}K_{p,s,\eta}
 \quad(1\le j\le s).
\]

\begin{proposition}[Integral Rees--Bol identity]\label{prop:Rees-Bol}
For $0\le j\le s$ there are integers $c_{j,h}$ such that
\begin{equation}\label{eq:Rees-Bol}
 \calK_j=\sum_{h=0}^jc_{j,h}X^hD^{j-h}K_{p,s,\eta},
 \qquad c_{j,0}=1.
\end{equation}
The transition from
$(K,DK,\ldots,D^{s-1}K)$ to
$(\calK_0,\ldots,\calK_{s-1})$ is lower triangular over $\Z[X]$
with diagonal $1$, and
\[
 \calK_s=D^sK=E_{p,s+1}^{\mathrm{evil}}.
\]
For $0\le j\le s-1$, the coefficient of $X^jK$ in $\calK_j$ is
\begin{equation}\label{eq:top-cj}
 \boxed{c_{j,j}=(-1)^j\frac{(s-1)!}{(s-1-j)!}.}
\end{equation}
The terminal case $j=s$ is the separate Bol identity
$\calK_s=D^sK$ and has no $X^sK$ term.
\end{proposition}

\begin{proof}
Induct using $\delta_k=D+kX$ and $D(X^h)=-hX^{h+1}$.  All coefficients
remain integral and the coefficient of the highest derivative remains one.
At the terminal step every $X$-term has a vanishing falling-factorial
coefficient, which is Bol's identity.  For the top $X$-coefficient write
$c_j=c_{j,j}$.  The next operator gives
$c_j=(j-s)c_{j-1}$ with $c_0=1$, proving \eqref{eq:top-cj}.
\end{proof}

\subsection{Functional rank and the zero estimate}

\begin{proposition}[Functional rank]\label{prop:functional-rank}
The complex monodromy orbit of the homogenized horizontal leaf
$S_{p,s,\eta}$ spans a space of dimension $s+1$.
\end{proposition}

\begin{proof}
The period-polynomial cocycle is not identically zero.  Otherwise, subtracting a polynomial
of degree at most $s-1$ from the Eichler primitive would produce a modular
form of weight $1-s<0$ with moderate growth at every cusp, hence the zero
form, contradicting $D^sK=E_{p,s+1}^{\mathrm{evil}}\ne0$.

Let $W$ be the span of the cocycle values in
$\Sym^{s-1}(\C^2)$.  The cocycle relation shows that $W$ is
$\Gamma_0(p)$-stable.  Since $\Gamma_0(p)$ is Zariski dense in
$\mathrm{SL}_2(\C)$ and $\Sym^{s-1}(\C^2)$ is irreducible, $W$ is the
full $s$-dimensional representation.  Adjoining the homogenizing constant
gives dimension $s+1$.
\end{proof}

\begin{lemma}[Functional independence]\label{lem:functional-independence}
In any rational algebraic frame of $\calE_{p,s}$, the coordinate functions of
the horizontal leaf are linearly independent over $\C(f)$, and hence over
$\Q(f)$.
\end{lemma}

\begin{proof}
If a $\C(f)$-valued rational covector annihilates one branch, its
single-valuedness implies that it annihilates every analytic continuation of
that branch.  At a generic basepoint it therefore annihilates the monodromy
orbit, which spans the fibre by \cref{prop:functional-rank}.  The covector is
zero.
\end{proof}

\begin{corollary}[Bost pivot set]\label{cor:zero-estimate}
Let $V_D$ be the set of nonzero jumps of the vanishing-order filtration of the
degree-$D$ evaluation map defined in \cref{sec:sources}.  For every
$\varepsilon>0$,
\begin{equation}\label{eq:zero-estimate}
 \#V_D=(s+1)D+O_{p,s}(1),
 \qquad
 V_D\subset
 [0,(s+1+\varepsilon)D+O_{p,s,\varepsilon}(1)].
\end{equation}
\end{corollary}

\begin{proof}
Put $m=s+1$.  Choose one fixed rational algebraic frame and multiply by a
fixed rational normalizer which is regular and nonzero at $P$.  Since the
bundle and normalizer are fixed, the degree-$D$ source embeds into
\[
 \bigl(\Q[f]_{\le D+C}\bigr)^m
\]
for one constant $C$, and its evaluation is obtained from $m$ holonomic
power series which are independent over $\C(f)$ by
\cref{lem:functional-independence}.  The Chudnovsky--Osgood estimate in the
form \cite[Theorem~3.2.13]{CDT2024}, applied with polynomial degree $D+C$,
therefore bounds the vanishing order of every nonzero source element by
\[
 (m+\varepsilon)(D+C+1)+C_\varepsilon
 =(m+\varepsilon)D+O_{p,s,\varepsilon}(1),
\]
after replacing $\varepsilon$ by a slightly smaller positive number if
necessary.  This proves the asserted containment of $V_D$.

The evaluation map is injective: a source in its kernel would give a
$\C(f)$-linear relation among the coordinate functions.  Since every
nonzero graded quotient of the one-variable vanishing filtration is
one-dimensional, the number of jumps is the source rank.  By
\cref{lem:graded-source}, this is $mD+O_{p,s}(1)$.
\end{proof}

\section{Global sources and fixed-data bridges}\label{sec:sources}

Let
\[
 0=\calF^{s+1}\subset\calF^s\subset\cdots\subset\calF^0=\calE_{p,s}
\]
be the cyclic oper filtration, with rank-one graded lines $L_i$.  On the dual
bundle put
\[
 \calA_i=\Ann(\calF^{i+1})\subset\calE_{p,s}^\vee,
 \qquad0\le i\le s.
\]
Choose a fixed divisor $B$, disjoint from the Tate cusp $P$, clearing all
fixed poles used below, and define
\begin{equation}\label{eq:source-MD}
 M_D=H^0\left(X_0(p),\calE_{p,s}^\vee(DP+B)\right),
 \qquad
 M_{D,i}=H^0\left(X_0(p),\calA_i(DP+B)\right).
\end{equation}
Because the extension bundle is split as
$\calE_{p,s}=\calO\oplus V$, put also
\begin{equation}\label{eq:source-split}
 M_D^0=H^0\left(X_0(p),\calO(DP+B)\right),
 \qquad
 M_D^+=H^0\left(X_0(p),V^\vee(DP+B)\right).
\end{equation}
Thus $M_D=M_D^0\oplus M_D^+$, and Riemann--Roch gives
\begin{equation}\label{eq:positive-source-rank}
 \dim_\Q M_D^+=sD+O_{p,s}(1).
\end{equation}
After spreading out, the same direct-sum decomposition holds for every
$\Z_\ell$-source lattice outside the fixed model exceptional set.

\begin{lemma}[Exact graded source and saturation]\label{lem:graded-source}
For $D\gg_{p,s}1$,
\[
 M_{D,i}/M_{D,i-1}
 \simeq H^0\left(X_0(p),L_i^\vee(DP+B)\right).
\]
In particular,
\[
 \dim_\Q M_D=(s+1)D+O_{p,s}(1),
 \qquad
 \dim_\Q(M_{D,i}/M_{D,i-1})=D+O_{p,s}(1).
\]
After enlarging one fixed exceptional set $\Sigma$, the induced filtration of
the $\Z_\ell$-source lattice is saturated for every $\ell\notin\Sigma$,
uniformly in $D$.
\end{lemma}

\begin{proof}
Spread the fixed exact sequences
$0\to\calA_{i-1}\to\calA_i\to L_i^\vee\to0$, the point $P$, and the
divisor $B$ over $R=\Z[\Sigma^{-1}]$.  Relative Serre vanishing gives a
single $D_0$ such that the first higher direct image vanishes for all
$D\ge D_0$.  Pushforward is then exact and compatible with base change, and
the integral quotients are finite locally free.  Riemann--Roch on
$X_0(p)\simeq\mathbf P^1$ gives the dimensions.
\end{proof}

\begin{lemma}[Fixed neat-level descent and uniform spreading out]
\label{lem:neat-descent}
After enlarging $B$ and a fixed finite set $\Sigma$, the following hold.

\begin{enumerate}[label=\textup{(\roman*)}]
\item The de Rham frame torsor and Gauss--Manin connection may be constructed
on one fixed neat tame cover and descended, through their canonical descent
data, to the modular stack; the associated bundle
$\calO\oplus\Sym^{s-1}\calH$, its connection, canonical BGG lift, and
contraction pairing descend to $X_0(p)\setminus\abs B$.

\item For every $D\gg1$ and every $\lambda\in M_D$, the normalized scalar
pairing may be computed on a flat algebraic torsor chart and has the form
\begin{equation}\label{eq:torsor-decomp-bridge}
 f^D\langle\lambda,S\rangle
 =H_\lambda(f,u)+\sum_{e=1}^sQ_{e,\lambda}(f,u)R_e(f),
\end{equation}
where $u$ denotes fibre coordinates,
\[
 \operatorname{supp}_fQ_{e,\lambda}\subset[-C,D+C],
\]
and the fibre-variable degree is bounded by a constant depending only on
$p,s$.

\item All these objects have models over $\Z[\Sigma^{-1}]$, with $\Sigma$
independent of $D$.  Fixed changes of model or of $B$ contribute
$o(D^2)$ to global arithmetic degrees and summed pivot heights.
\end{enumerate}
\end{lemma}

\begin{proof}
Choose a fixed neat tame level $N\ge5$, prime to $p$, and a finite Galois
cover $Y\to\mathscr X_0(p)$.  The frame torsor and connection of
\cite[Propositions~A.14--A.15]{GPRJ2025} are functorial in the elliptic curve
and level structure, so the two pullbacks to
$Y\times_{\mathscr X_0(p)}Y$ agree and satisfy the cocycle condition.  They
descend to the modular stack.  The generic stabilizer $-1$ acts trivially on
$\Sym^{s-1}\calH$ because $s-1$ is even.  Put the finitely many elliptic
points into $B$; the associated bundle, the canonical lift of
\cref{lem:BGG}, and the invariant pairing then descend to the coarse curve
away from $B$.

Choose a fixed rational algebraic frame on that affine open, regular at $P$.
Since $f$ is a rational coordinate with a simple zero at $P$, multiplication
by $f^D$ cancels the variable pole at $P$.  The remaining coordinates have
$f$-support $[-C,D+C]$ and only fixed poles away from $P$.  Pullback to the
frame torsor and use the fixed representation $1\oplus\Sym^{s-1}$ together
with the fixed-order BGG/Rees--Bol differential operator.  This gives
\eqref{eq:torsor-decomp-bridge} with fibre degree bounded solely in terms of
$p,s$.

Finally, the finite collection of bundles, matrices, divisors, torsor charts,
and morphisms spreads out after clearing finitely many denominators.  The
torsor chart may be chosen smooth, hence flat, over the resulting base.  A fixed model change distorts norms on a source of rank $O(D)$ and on
$O(D)$ one-dimensional pivot quotients, hence contributes at most
$O(D\log D)=o(D^2)$.
\end{proof}

\begin{proposition}[Single unipotent action on the deepest multiplier]
\label{prop:single-unipotent}
Write $s=2r+1$ and $H=D\log f$.  Fix the algebraic Hodge splitting
$\sigma$ used below, regular at $P$, and write $x=X/H$ for the normalized
unipotent coordinate.  There are a fixed
polynomial $g(f)\in\Z[\Sigma^{-1}][f]$ with $g(0)$ a unit and a constant $C$
such that, for every $D\gg1$ and every $\lambda\in M_D^+$, the multiplier of
the normalized undifferentiated row $H^rK$ in
$f^D\langle\lambda,S\rangle$ has the form
\begin{equation}\label{eq:genuine-q0-shape}
 q_{0,\lambda}(f,x)=\sum_{j=0}^{s-1}Q_{j,\lambda}(f)x^j,
\end{equation}
linear in $\lambda$, with
\begin{equation}\label{eq:genuine-q0-degree}
 g(f)Q_{j,\lambda}(f)\in\Z[\Sigma^{-1}][f],
 \qquad
 \deg\bigl(gQ_{j,\lambda}\bigr)\le D+C.
\end{equation}
The same assertions hold after reduction at every $\ell\notin\Sigma$.
Only one degree-$(s-1)$ unipotent matrix occurs in
\eqref{eq:genuine-q0-shape}; in particular there is no doubled
$2(s-1)$ fibre degree.
\end{proposition}

\begin{proof}
Put $n=s-1$ and choose the monomial basis
$e_j=e_0^{n-j}e_1^j$ of $\Sym^n\calH$.  For
\[
 u(x)=\begin{pmatrix}1&x\\0&1\end{pmatrix}
\]
one has, in this basis,
\begin{equation}\label{eq:unipotent-matrix}
 \Sym^n(u(x))e_j
 =\sum_{h=0}^j\binom{j}{h}x^{j-h}e_h.
\end{equation}
Thus every entry of $\Sym^n(u(x))$, its inverse, and their dual matrices is
an integral polynomial in $x$ of degree at most $n$.

Write the coordinates of $f^D\lambda$ in the dual frame induced by the fixed
splitting $\sigma$ as
$P_{0,\lambda}(f),\ldots,P_{n,\lambda}(f)$.  These are functions on the base
and contain no fibre coordinate.  In the canonical BGG lift
\eqref{eq:spl-def}, the undifferentiated $K$ occurs only in the bottom
extremal component.  The Tate oper frame and the $\sigma$-frame differ by a
diagonal torus factor and one unipotent factor $u(x)$.  Transport the
$V$-valued BGG lift to the $\sigma$-frame and keep the source covector in the
$\sigma$-dual frame.  By equivariance of contraction,
\[
 \langle P,\rho(u(x))W\rangle
 =\langle\rho^\vee(u(x)^{-1})P,W\rangle,
 \qquad \rho=\Sym^n,
\]
so either description uses one matrix from \eqref{eq:unipotent-matrix}, not a
product of two such matrices.  More explicitly, let $\operatorname{coeff}_{K}$ denote projection onto
the literal undifferentiated raw $K$-row and normalize the coefficient by
$H^{-r}$.  Then
\begin{align}
 q_{0,\lambda}(f,x)
 &=H^{-r}\operatorname{coeff}_{K}
   \left\langle f^D\lambda,
     \rho(u(x))\operatorname{spl}_s(K)\right\rangle\label{eq:q0-one-action}\\
 &=H^{-r}\operatorname{coeff}_{K}
   \left\langle\rho^\vee(u(x)^{-1})f^D\lambda,
     \operatorname{spl}_s(K)\right\rangle.\nonumber
\end{align}
The Rees--Bol identity is the raw-derivative coordinate formula for this same
transported BGG section; it is not a second frame transport.  Since the
undifferentiated $K$ occurs only in the bottom extremal component of
\eqref{eq:spl-def}, \eqref{eq:q0-one-action} is a fixed row of one matrix in
\eqref{eq:unipotent-matrix} applied to the base-coordinate vector
$P_\lambda$.  Hence it has $x$-degree at most $s-1$.

The factor $H^{-r}$ is independent of $x$, and it exactly removes the
diagonal weight of the bottom component.  The resulting coefficient is
weight zero and therefore descends from the torsor to a rational function on
the base.  More generally the diagonal factor contributes no additional
variable because
\[
 \operatorname{wt}(\calK_j)=1-s+2j,
 \qquad \operatorname{wt}(H)=2,
 \qquad \operatorname{wt}(H^{r-j}\calK_j)=0.
\]
Finally, as in \cref{lem:neat-descent}, a fixed polynomial $g(f)$, nonzero at
$P$, clears all fixed poles of the coordinates $P_{j,\lambda}$ and the
fixed poles introduced by the normalization $H^{-r}$; multiplication by
$f^D$ then gives degree $D+O(1)$.  This proves
\eqref{eq:genuine-q0-shape}--\eqref{eq:genuine-q0-degree}.  The entire
construction spreads out with the fixed torsor and representation, so the
same formula holds after every reduction at $\ell\notin\Sigma$.
\end{proof}

\begin{lemma}[Adelic fixed-bundle norm comparison]
\label{lem:fixed-bundle}
Let $\calV$ be a fixed rank-$m$ vector bundle on $\mathbf P^1_\Q$, equipped
with fixed adelic models and metrics, and put
\[
 N_D=H^0\left(\mathbf P^1,\calV(DP+B)\right).
\]
Suppose that the scalar contraction with a fixed horizontal leaf is
injective and that its pivot set satisfies a linear zero estimate.  Relative
to $m$ scalar polynomial sources of degree $D$, replacing the source by
$N_D$, shifting the scalar degrees by fixed integers, multiplying by a fixed
normalizer regular and nonzero at $P$, changing a fixed algebraic frame, or
changing finitely many local models alters both the source degree and the
sum of evaluation heights by $o(D^2)$.  Consequently
\begin{equation}\label{eq:fixed-bundle-degree}
 \widehatdeg N_D=\frac m2\Energy D^2+o(D^2),
\end{equation}
and
\begin{equation}\label{eq:fixed-bundle-heights}
 \sum_{n\in V_D}\sum_vh_v(\psi_D^{(n)})
 \le\left(-\frac{m^2}{2}\LambdaAH+m\Energy\right)D^2
 +\mathscr A_D+o(D^2).
\end{equation}
\end{lemma}

\begin{proof}
By Birkhoff--Grothendieck there is a fixed global bundle isomorphism
\begin{equation}\label{eq:BG-splitting}
 \calV(B)\simeq\bigoplus_{i=1}^m\calO(a_i)
\end{equation}
with fixed integers $a_i$.  Transport the fixed adelic data through this
isomorphism.  At every place, the original and transported norms are
uniformly equivalent on the fixed analytic template; outside a finite set of
finite places the integral models agree up to units.  Thus there are constants
$C_v\ge1$, equal to $1$ for all but finitely many $v$, which bound the
operator norms of the comparison and its inverse.  If
$N_D^{\mathrm{spl}}$ denotes the same source with the transported split norms,
then
\[
 \left|\widehatdeg N_D-\widehatdeg N_D^{\mathrm{spl}}\right|
 \le \rank(N_D)\sum_v\log C_v=O(D).
\]
This is the required determinant-norm comparison.

Under \eqref{eq:BG-splitting}, the source is the direct sum of the scalar
sources
\[
 H^0\bigl(\mathbf P^1,\calO((D+a_i)P)\bigr),
 \qquad 1\le i\le m,
\]
up to fixed conventions for the divisor $B$.  The scalar source-degree asymptotic, in the normalization of
\cite[(7.3.7)]{CDT2024}, and its adelic form in
\cite[Section~2.3 and the proof of Theorem~2.6]{CDT2025}, where the
arithmetic Hilbert--Samuel theorem is combined with the Bost--Charles overflow
identities \cite[Theorem~5.4.1 and Proposition~5.4.2]{BostCharles2022}, gives
\[
 \frac12\Energy(D+a_i)^2
 =\frac12\Energy D^2+O(D)
\]
for the $i$th summand.  Summing and including the $O(D)$ norm distortion
proves \eqref{eq:fixed-bundle-degree}.

We next compare evaluation heights.  In the split frame the contraction is a
scalar Pad\'e evaluation with $m$ fixed holonomic coordinate functions and
row degrees $D+a_i$.  On each fixed continuation template the bundle
comparison matrices are analytic and bounded; after the fixed normalizer has
cleared the fixed poles, the same is true of the matrices needed for the
local upper bound.  At all but finitely many finite places these matrices are
integral of norm one.  The normalizer is regular and nonzero at $P$, so it
preserves vanishing order, and its local sup norm contributes only a fixed
factor.  Consequently these comparisons contribute $O(1)$ per pivot;
allowing the standard monomial/coefficient-lattice comparison gives the
harmless bound $O(\log D)$ per pivot.  There are $mD+O(1)$ pivots, so the total comparison error is
$O(D\log D)=o(D^2)$.

Replacing $D+a_i$ by $D$ changes the analytic $D\Energy$ term by $O(1)$ per
pivot.  It also shifts any finite-place multiplier window by only $O(1)$;
the sum of the corresponding fixed boundary contributions is
$O(D\log D)$.  The scalar CDT estimates
\cite[(7.3.14)--(7.3.15)]{CDT2024} therefore apply with unchanged $D^2$
coefficients.  Finally, the $mD+O(1)$ pivot orders are distinct and
nonnegative, whence
\[
 \sum_{n\in V_D}n\ge\frac12m^2D^2+o(D^2).
\]
Inserting this in the scalar height estimate proves
\eqref{eq:fixed-bundle-heights}.
\end{proof}

\section{The hybrid local arithmetic theorem}\label{sec:hybrid-local}

Throughout this section, $p$ and the odd integer $s$ are fixed, and
$\zeta_p(s)\in\Q$ is assumed.  Put
\[
 \eta=\zeta_p(s)/2\in\Q.
\]
Let $\Sigma_{p,s}$ be the fixed geometric/model exceptional set, independent
of $\eta$, enlarged to contain $p$ and every prime
$\ell\le\max\{3,s\}$.  Write $\eta=u/v$ in lowest terms, with
$u\in\Z$, $v\in\Z_{>0}$, and put
\[
 \Sigma_\eta=\{\ell:\ell\mid v\},
 \qquad t_\ell=v_\ell(v).
\]
We call a prime $\ell$ \emph{geometrically good} if
$\ell\notin\Sigma_{p,s}$, and \emph{arithmetically good} if it is
geometrically good and $\ell\notin\Sigma_\eta$.  Thus every geometrically
good prime automatically satisfies $\ell>\max\{3,s\}$.  The Hasse
codimension theorem below is uniform at every geometrically good prime and is
independent of $\eta$; the one-power saving is used only at arithmetically
good primes.  The explicit baseline lattice at the finite set
$\Sigma_{p,s}\cup\Sigma_\eta$ is defined after the shifted cutoff $N^\sharp$ is fixed in
\eqref{eq:bad-small-lattice}.  Constants below are uniform in $D$ and in every
prime to which the corresponding statement applies.

\subsection{Weight-zero normalization and the comparison coordinate}

Write $s=2r+1$ and put
\begin{equation}\label{eq:A-C-rows}
 A_j=H^{r-j}D^jJ,
 \qquad
 C_j=H^{r-j}\calK_j,
 \qquad
 H=D\log f.
\end{equation}
Since $H=1+O(q)$, its positive and negative fixed powers are integral formal
units outside a fixed exceptional set.  With $x=X/H$, \cref{prop:Rees-Bol}
and $K=J+\eta$ give, for $0\le j\le s-1$,
\begin{equation}\label{eq:C-vs-A}
 \boxed{
 C_j=\sum_{h=0}^jc_{j,h}x^hA_{j-h}
      +c_{j,j}\eta H^r x^j.}
\end{equation}
Indeed $D^{j-h}K=D^{j-h}J$ unless $h=j$, when the undifferentiated
constant $\eta$ remains.  Thus, for a formal linear combination of these
Rees--Bol iterates with multipliers $P_j$, if
\begin{equation}\label{eq:qk-definition}
 q_k=\sum_{j=k}^{s-1}c_{j,j-k}x^{j-k}P_j
 \qquad(0\le k\le s-1),
\end{equation}
then
\begin{equation}\label{eq:corrected-row-decomp}
 \boxed{
 \sum_{j=0}^{s-1}P_jC_j
 =q_0(A_0+\eta H^r)+\sum_{k=1}^{s-1}q_kA_k.}
\end{equation}
In particular, the rational constant term and the depth-$s$ row have the same
multiplier $q_0$.  This is an identity in the polynomial realization of the
BGG lift; it is not an assertion that the functions $C_j$ form a rational
scalar coordinate system for the global source.  Actual global contractions
are handled by \cref{eq:torsor-decomp-bridge,prop:single-unipotent}.

Choose an algebraic Hodge splitting $\sigma$ on an affine open containing the
Tate cusp.  Evaluating Urban's operator on a classical weight-$k$ modular
form gives $D_q+kX_\sigma$, whereas the Serre derivative is
$D_q-kE_2/12$.  Their difference is a meromorphic modular form of weight two.
Thus
\begin{equation}\label{eq:X-sigma}
 X_\sigma=-\frac{E_2}{12}+g_\sigma,
 \qquad g_\sigma\in M_2^{\mathrm{mer}}(\Gamma_0(p)),
\end{equation}
and hence
\begin{equation}\label{eq:x-sigma}
 x_\sigma=-\frac{E_2}{12H}+r_\sigma(f),
 \qquad r_\sigma(f)\in\Q(f).
\end{equation}

\begin{lemma}[Hasse algebraization and degree]\label{lem:Hasse-degree}
For every geometrically good auxiliary prime $\ell$, the reduction of
$x_\sigma$ belongs to $\F_\ell(f)$.  Writing
\[
 \bar x_\sigma=\frac{a_\ell(f)}{b_\ell(f)}
\]
in lowest terms with $b_\ell(0)\ne0$, one has
\begin{equation}\label{eq:delta-bound}
 \delta_\ell:=\max\{\deg a_\ell,\deg b_\ell\}
 \le c_p\ell+C,
 \qquad c_p=\frac{p+1}{12}.
\end{equation}
\end{lemma}

\begin{proof}
Modulo $\ell$, the Hasse invariant is represented by $E_{\ell-1}$, whose
$q$-expansion is $1$, and $E_{\ell+1}\equiv E_2$.  Hence the reduction of
\eqref{eq:x-sigma} is represented by
\[
 -\frac{E_{\ell+1}}{12H E_{\ell-1}}+\bar r_\sigma(f),
\]
a modular function on $X_0(p)_{\F_\ell}$.  The numerator and denominator of
the first term are sections of the same weight-$(\ell+1)$ line bundle.  The
classical weighted valence formula for $\Gamma_0(p)$ gives total orbifold zero
degree
\[
 \frac{[\mathrm{SL}_2(\Z):\Gamma_0(p)](\ell+1)}{12}
 =\frac{(p+1)(\ell+1)}{12}
\]
for either section; see \cite{Katz1973}.  We use only this upper bound and its
leading coefficient.  Since the numerator and denominator have the same
weight, their quotient is weight zero.  On passage to the coarse genus-zero
curve, the only convention change occurs at the finitely many elliptic
points; the correction to the zero or pole order is bounded independently of
$\ell$.  Adding the fixed rational term $\bar r_\sigma(f)$ changes numerator
and denominator degrees by another $O_{p,s}(1)$.  Hence
\[
 \max\{\deg a_\ell,\deg b_\ell\}
 \le\frac{p+1}{12}\ell+O_{p,s}(1),
\]
which is \eqref{eq:delta-bound}.  Regularity of the chosen splitting at the
cusp and $H(0)=1$ give $b_\ell(0)\ne0$.
\end{proof}

\subsection{The genuine-source Hasse kernel}

The small-prime argument is made on the actual global positive source
$M_D^+$, not on a putative scalar Rees--Bol coordinate block on the base.
For $\lambda\in M_D^+$ let $q_{0,\lambda}(f,x)$ be the multiplier of
the normalized deepest row $H^rK$ from \cref{prop:single-unipotent}; after
separating $K=J+\eta$ it multiplies $A_0+\eta H^r$.  On the splitting $\sigma$ it is evaluated at
$x=x_\sigma(f)$.

\begin{theorem}[Uniform genuine-source Hasse codimension]
\label{thm:genuine-Hasse}
There are constants $D_0,C$, depending only on $p,s$ and the fixed models,
such that the following holds.  Let $D\ge D_0$ and let
$\ell$ be geometrically good.  Write
\[
 \bar x_\sigma=\frac{a_\ell(f)}{b_\ell(f)},
 \qquad
 \delta_\ell=\max\{\deg a_\ell,\deg b_\ell\}
\]
as in \cref{lem:Hasse-degree}.  Then the reduction of the genuine positive
source contains a subspace
$\overline K_{D,\ell}\subset\overline M_D^+$ with
\begin{equation}\label{eq:kernel-rank}
 \boxed{
 \dim_{\F_\ell}\overline K_{D,\ell}
 \ge (s-1)\max\{0,D-\delta_\ell-C\}
 \ge (s-1)\max\{0,D-\ceil{c_p\ell}-C\},}
\end{equation}
such that
\begin{equation}\label{eq:q0-mod-l-zero}
 q_{0,\bar\lambda}\!\left(f,\frac{a_\ell}{b_\ell}\right)=0
 \quad\text{in }\F_\ell[[f]]
 \qquad(\bar\lambda\in\overline K_{D,\ell}).
\end{equation}
Any subspace of $\overline K_{D,\ell}$ lifts to a saturated direct summand
$K_{D,\ell}\subset M^+_{D,\ell}$ satisfying
\begin{equation}\label{eq:q0-l-divisible}
 q_{0,\lambda}(f,x_\sigma(f))\in\ell\Z_\ell[[f]]
 \qquad(\lambda\in K_{D,\ell}).
\end{equation}
The constants, the geometric exceptional set, and this codimension statement
are independent of $\eta$.
\end{theorem}

\begin{proof}
By \cref{prop:single-unipotent}, after multiplication by the fixed polynomial
$g(f)$ with $g(0)$ a unit, one has
\[
 q_{0,\lambda}(f,x)=\sum_{j=0}^{s-1}Q_{j,\lambda}(f)x^j,
 \qquad
 \deg(gQ_{j,\lambda})\le D+C_0.
\]
Reduction modulo $\ell$ therefore defines an $\F_\ell$-linear map on the
actual global source
\begin{align}
 T_{D,\ell}:\overline M_D^+&\longrightarrow
 \F_\ell[f]_{\le D+(s-1)\delta_\ell+C_1},\label{eq:Hasse-T-map}\\
 \bar\lambda&\longmapsto
 \bar g(f)b_\ell(f)^{s-1}
 q_{0,\bar\lambda}\!\left(f,\frac{a_\ell(f)}{b_\ell(f)}\right).
 \nonumber
\end{align}
Indeed, the $j$th summand after clearing $b_\ell^{s-1}$ has degree at most
$D+C_0+j\delta_\ell+(s-1-j)\delta_\ell$.

By \eqref{eq:positive-source-rank} and base change,
\[
 \dim_{\F_\ell}\overline M_D^+=sD+O_{p,s}(1),
\]
whereas the target of \eqref{eq:Hasse-T-map} has dimension at most
$D+(s-1)\delta_\ell+O_{p,s}(1)$.  Rank--nullity gives
\[
 \dim\ker T_{D,\ell}
 \ge(s-1)(D-\delta_\ell)-C_2.
\]
After increasing one fixed constant $C$, this implies the first lower bound
in \eqref{eq:kernel-rank}; \cref{lem:Hasse-degree} gives the second.  Take
$\overline K_{D,\ell}$ to be a subspace of the kernel of the displayed
conservative dimension.

Since $g(0)$ and $b_\ell(0)$ are units, both are invertible in
$\F_\ell[[f]]$.  The identity $T_{D,\ell}(\bar\lambda)=0$ therefore implies
\eqref{eq:q0-mod-l-zero}; equivalently, use that
$\F_\ell[[f]]$ is a domain.  Extend a basis of
$\overline K_{D,\ell}$ to an $\F_\ell$-basis of
$\overline M_D^+$ and lift the entire basis to the free lattice
$M_{D,\ell}^+$.  The lifted first vectors span a saturated direct summand
$K_{D,\ell}$.  The integral models of \cref{lem:neat-descent} make the
multiplier integral and its reduction compatible with this lift.  Its
reduction is zero, proving \eqref{eq:q0-l-divisible}.
\end{proof}

\begin{remark}
The proof uses only the actual global source, the degree-$(s-1)$ unipotent
representation, and Hasse algebraization of $x_\sigma$.  It does not identify
the normalized Rees--Bol rows with scalar coordinates in a rational frame on
the base curve.
\end{remark}

\begin{lemma}[Actual deepest-row decomposition]\label{lem:actual-deepest-row}
On the fixed splitting chart, for every $D\gg1$ and every
$\lambda\in M_D^+$ one may write the normalized genuine contraction as
\begin{equation}\label{eq:actual-deepest-row}
 f^D\langle\lambda,S\rangle
 =G_\lambda(f,u)+q_{0,\lambda}(f,x_\sigma)H^rJ
  +\eta q_{0,\lambda}(f,x_\sigma)H^r
  +\sum_{e=1}^{s-1}\widetilde Q_{e,\lambda}(f,u)R_e(f),
\end{equation}
where $G_\lambda$ is depth zero, the $f$-supports are contained in a fixed
enlargement of $[-C,D+C]$, and all displayed multipliers are integral at every
arithmetically good prime.  Thus the multiplier of the unique depth-$s$ row
$R_s=J$ is the formal unit $H^r$ times the genuine multiplier
$q_{0,\lambda}$ of \cref{prop:single-unipotent}; every other nonalgebraic row
has depth at most $s-1$.
\end{lemma}

\begin{proof}
Apply the genuine torsor decomposition \eqref{eq:torsor-decomp-bridge} to the
canonical BGG lift.  By \cref{prop:single-unipotent}, the coefficient of its normalized
undifferentiated component $H^rK=H^r(J+\eta)$ is $q_{0,\lambda}$.
Separating $K=J+\eta$ gives the two middle terms of
\eqref{eq:actual-deepest-row}.  All remaining components involve
$D^jJ=R_{s-j}$ with $j\ge1$, hence have depth at most $s-1$.  The support and
integrality assertions follow from the spread-out model in
\cref{lem:neat-descent}; fixed powers of $H$ are integral formal units at an
arithmetically good prime.
\end{proof}

\subsection{The small-prime hybrid lattice}

Fix $1<\xi<s+1$.  Choose once and for all a constant
$C_{\mathrm{sh}}$ at least as large as the absolute value of every fixed
negative $f$-support bound in \cref{lem:neat-descent,lem:actual-deepest-row},
and put
\[
 N_0=\floor{\xi D},
 \qquad
 N^\sharp=N_0+C_{\mathrm{sh}},
 \qquad
 a_{\ell,D}=v_\ell(L_{N^\sharp}).
\]
The fixed enlargement from $N_0$ to $N^\sharp$ ensures that a Laurent
multiplier with lower support at least $-C_{\mathrm{sh}}$ can only use raw
coefficients of index at most $N^\sharp$ in an evaluated coefficient through
$f^{N_0}$.  Since $C_{\mathrm{sh}}$ is fixed,
$\log L_{N^\sharp}=\xi D+o(D)$.
For every $\ell\in\Sigma_{p,s}$ choose a fixed integer
$\beta_\ell\ge0$ clearing the denominators of the chosen local model, frame,
and fixed powers of $H$; put $\beta_\ell=0$ for
$\ell\notin\Sigma_{p,s}$.  At every prime
$\ell\in\Sigma_{p,s}\cup\Sigma_\eta$ with $\ell\le N^\sharp$, define the baseline
lattice
\begin{equation}\label{eq:bad-small-lattice}
 M^{\mathrm{bad}}_{D,\ell}
 =\ell^{\beta_\ell}M^0_{D,\ell}\oplus
 \ell^{sa_{\ell,D}+\beta_\ell+t_\ell}M^+_{D,\ell},
\end{equation}
where $t_\ell=0$ when $\ell\notin\Sigma_\eta$.  The factor
$\ell^{sa_{\ell,D}}$ clears every positive-depth raw row needed through
$f^{N_0}$ by \cref{lem:raw-depth}; the fixed factor $\beta_\ell$ clears
model and frame denominators on both split source blocks, and
$t_\ell$ clears the denominator of $\eta$.  Since the exceptional set is
finite and $a_{\ell,D}=O_\ell(\log D)$, its total index contribution is
$O_{p,s,\eta}(D\log D)=o(D^2)$.  Omitting the one-power Hasse saving at these
fixed primes costs only $O_{p,s,\eta}(D)$.

For each arithmetically good prime $\ell\le N^\sharp$, use
\cref{thm:genuine-Hasse} to choose a
saturated direct summand $K_{D,\ell}\subset M^+_{D,\ell}$ of conservative
rank
\begin{equation}\label{eq:r-star}
 r^*_{D,\ell}=(s-1)
 \max\{0,D-\ceil{c_p\ell}-C\},
\end{equation}
and choose a complement $K^c_{D,\ell}$ in $M^+_{D,\ell}$.  With the exact
split source $M_{D,\ell}=M^0_{D,\ell}\oplus M^+_{D,\ell}$, define
\begin{equation}\label{eq:hybrid-lattice}
 M^{\mathrm{hyb}}_{D,\ell}
 =M^0_{D,\ell}
 \oplus \ell^{sa_{\ell,D}-1}K_{D,\ell}
 \oplus \ell^{sa_{\ell,D}}K^c_{D,\ell}.
\end{equation}

\begin{lemma}[Small-prime integrality]\label{lem:small-integrality}
Every unit vector of \eqref{eq:hybrid-lattice} has integral evaluation through
$f^{N_0}$.
\end{lemma}

\begin{proof}
Write $a=a_{\ell,D}\ge1$.  For $\lambda\in K_{D,\ell}$ use the actual
genuine-source decomposition \eqref{eq:actual-deepest-row}.  By
\cref{thm:genuine-Hasse},
$q_{0,\lambda}(f,x_\sigma(f))\in\ell\Z_\ell[[f]]$, and $H^r$ is a formal
unit.  Hence the depth-$s$ row $J=R_s$ acquires the missing one power of
$\ell$, so $(sa-1)+1=sa$ clears its $L_{N^\sharp}^s$ denominator.  The constant term is
integral because $\eta\in\Z_\ell$ at an arithmetically good prime and is
multiplied by the same divisible coefficient.

Every remaining nonalgebraic row in \eqref{eq:actual-deepest-row} has depth
$d\le s-1$, while
\[
 sa-1\ge da
\]
for $a\ge1$.  Its multiplier is integral by
\cref{lem:actual-deepest-row}.  Thus $\ell^{sa-1}$ clears every row on
$K_{D,\ell}$; the complement is scaled by the worst depth $sa$.  Apply
\cref{lem:raw-depth} with cutoff $N^\sharp$ after the integral
substitution $q=q(f)$.
\end{proof}

\begin{proposition}[Small-prime determinant cost]\label{prop:small-cost}
The local index exponent is
\begin{equation}\label{eq:local-index}
 v_\ell[M_{D,\ell}:M^{\mathrm{hyb}}_{D,\ell}]
 =s^2Da_{\ell,D}-r^*_{D,\ell}+O_{p,s}(a_{\ell,D}).
\end{equation}
After summing over arithmetically good $\ell\le N^\sharp$, and using
\eqref{eq:bad-small-lattice} at the finite exceptional set, the total
small-prime index is
\begin{equation}\label{eq:small-cost}
 \sum_{\ell\le N^\sharp}
 \log[M_{D,\ell}:M^{\mathrm{small}}_{D,\ell}]
 =S_{p,s}(\xi)D^2+o(D^2),
\end{equation}
where $M^{\mathrm{small}}_{D,\ell}$ denotes
$M^{\mathrm{hyb}}_{D,\ell}$ at arithmetically good primes and
$M^{\mathrm{bad}}_{D,\ell}$ at primes in
$\Sigma_{p,s}\cup\Sigma_\eta$.  These cases are exhaustive by construction.
We define the adelic source lattice
$M_D^{\mathrm{hyb}}$ by these local lattices for $\ell\le N^\sharp$ and by the
unmodified source lattice for $\ell>N^\sharp$.  Here
\[
 S_{p,s}(\xi)=s^2\xi-(s-1)
 \int_0^\xi(1-c_px)_+\,\dd x.
\]
In particular, if $c_p\xi<1$, then
\begin{equation}\label{eq:small-cost-simple}
 S_{p,s}(\xi)=s^2\xi-(s-1)
 \left(\xi-\frac{p+1}{24}\xi^2\right).
\end{equation}
\end{proposition}

\begin{proof}
The positive source has rank $sD+O(1)$.  Hence
\[
 (sa-1)r^*+sa(sD+O(1)-r^*)
 =s^2Da-r^*+O(a),
\]
which proves \eqref{eq:local-index}.  Now
\[
 \sum_{\ell\le N^\sharp}a_{\ell,D}\log\ell
 =\log L_{N^\sharp}=\xi D+o(D).
\]
The prime number theorem and partial summation give
\[
 \sum_{\ell\le N^\sharp}r^*_{D,\ell}\log\ell
 =(s-1)D^2\int_0^\xi(1-c_px)_+\,\dd x+o(D^2).
\]
Removing the finite set $\Sigma_{p,s}\cup\Sigma_\eta$ from the two prime
sums changes their leading terms by at most $O_{p,s,\eta}(D\log D)$.  The
baseline lattices \eqref{eq:bad-small-lattice} and the $O(a)$ errors also
contribute $O_{p,s,\eta}(D\log D)=o(D^2)$.
\end{proof}

\begin{lemma}[Residual small-prime heights]\label{lem:small-residual}
After the cutoff, the total residual positive denominator penalty at primes
$\ell\le N^\sharp$, summed over all pivots, is
\[
 O_{p,s,\xi}(D^{3/2}\log D)=o(D^2).
\]
\end{lemma}

\begin{proof}
Every coefficient index occurring at a pivot is $O_s(D)$ by
\cref{cor:zero-estimate}.  A residual prime power can occur only if
\[
 \ell^{a_{\ell,D}+1}>N^\sharp,
 \qquad
 \ell^{a_{\ell,D}+1}=O_s(D).
\]
Since $a_{\ell,D}+1\ge2$, this forces $\ell=O_s(\sqrt D)$.  Summing proper
prime-power excesses gives $O_s(\sqrt D\log D)$ for one pivot and therefore
$O_s(D^{3/2}\log D)$ over $(s+1)D+O(1)$ pivots.  This includes the finite bad
primes: the fixed factors $\beta_\ell+t_\ell$ in
\eqref{eq:bad-small-lattice} clear their model and constant denominators, and
only the same prime-power excess beyond $L_{N^\sharp}$ remains.  The one omitted Hasse
power at arithmetically good primes is exactly restored by the deepest-row
divisibility and creates no additional residual term.
\end{proof}

\subsection{Large primes on the de Rham torsor}

At primes $\ell>N^\sharp$ use the unmodified source lattice.  By
\cref{lem:neat-descent}, after pullback to a flat integral torsor chart the
normalized scalar evaluation has the decomposition
\begin{equation}\label{eq:torsor-eval}
 F_\lambda(f)=H_\lambda(f,u)+
 \sum_{e=1}^sQ_{e,\lambda}(f,u)R_e(f),
\end{equation}
with $f$-support $[-C,D+C]$ and fixed fibre degree.  The value of
$F_\lambda$ is independent of the torsor point because it is the contraction
of a source with a horizontal section.

Insert \cref{lem:Dwork} into \eqref{eq:torsor-eval}:
\begin{equation}\label{eq:torsor-Dwork}
 F_\lambda=H_\lambda^\flat+
 \sum_{e=1}^s\ell^{-e}Q_{e,\lambda}(f,u)
 R_e(\varphi_\ell(f)),
\end{equation}
where $H_\lambda^\flat$ is integral.  The fibre variables are coefficients;
no Frobenius action on them is used.

\begin{lemma}[Coefficient-ring Taylor no-backflow]\label{lem:no-backflow}
Let $A$ be any $\F_\ell$-algebra and
\[
 \Gamma(f,z)=\sum_{j=0}^df^jP_j(z),
 \qquad d<\ell,
 \qquad P_j(z)\in A[[z]].
\]
Suppose $\Gamma(f,f^\ell)\in f^nA[[f]]$ and
$\Delta(f)\in f^{\ell+1}A[[f]]$.  Then for every $t\ge0$,
\begin{equation}\label{eq:no-backflow}
 \Delta(f)^t
 \left.\partial_z^t\Gamma(f,z)\right|_{z=f^\ell}
 \in f^{n+t}A[[f]].
\end{equation}
\end{lemma}

\begin{proof}
Write $P_j(z)=\sum_rp_{j,r}z^r$.  Since $0\le j<\ell$, the exponent
$j+\ell r$ uniquely determines $(j,r)$; no cancellation between different
pairs is possible, even if $A$ has zero divisors.  Thus a nonzero
$p_{j,r}$ satisfies $j+\ell r\ge n$.  Terms with $r<t$ vanish after $t$
derivatives.  For $r\ge t$, after differentiating $t$ times and multiplying
by $\Delta^t$, the corresponding order is at least
\[
 j+\ell(r-t)+t(\ell+1)=j+\ell r+t\ge n+t.
\]
\end{proof}

\begin{lemma}[Pivot-range truncation]\label{lem:pivot-truncation}
Let $\ell>N^\sharp$ and let
\[
 F_\lambda(f)=c_nf^n+O(f^{n+1}),
 \qquad c_n\ne0,
\]
be a Bost pivot.  Put
$\widetilde F_\lambda=f^CF_\lambda$ and $\widetilde n=n+C$, where the fixed
$C$ shifts all multiplier $f$-degrees into $[0,D+2C]$.  For all sufficiently
large $D$, every coefficient through order $f^{\widetilde n}$ in
\eqref{eq:torsor-Dwork} is unchanged if each $R_e(z)$ is replaced by its
truncation $R_{e,<\ell}(z)$ below $z^\ell$.  The coefficients of every
$R_{e,<\ell}$ are $\ell$-integral.  Consequently, if
$v_\ell(c_n)<0$, then
\begin{equation}\label{eq:pivot-pole-bound}
 1\le -v_\ell(c_n)\le s.
\end{equation}
\end{lemma}

\begin{proof}
By \cref{cor:zero-estimate}, $\widetilde n=O_{p,s}(D)$.  Since
$\ell>N^\sharp=\floor{\xi D}+C_{\mathrm{sh}}$ with $\xi>1$, one has
$D+2C<\ell$ and $\widetilde n<\ell^2$ for all sufficiently large $D$.
Consider a monomial $f^jz^r$ in a multiplier times $R_e(z)$.  For a
nonzero $t$th Taylor term one has $t\le r$, and at
$z=f^\ell+\ell\Delta_\ell(f)$ its $f$-order is at least
\[
 j+\ell(r-t)+t(\ell+1)=j+\ell r+t.
\]
If $r\ge\ell$, this is at least $\ell^2$, and hence is strictly beyond the
pivot order.  Thus terms of $z$-degree at least $\ell$ cannot affect
$c_n$.  For $r<\ell$, the index $r$ is an $\ell$-adic unit, so the raw-depth
formula makes the coefficient of $z^r$ in $R_e(z)$ $\ell$-integral.  After
truncation the only negative power of $\ell$ in
\eqref{eq:torsor-Dwork} is therefore the explicit factor $\ell^{-e}$, with
$1\le e\le s$.  This proves \eqref{eq:pivot-pole-bound}.
\end{proof}

Let $A_\ell$ be the integral coefficient algebra of the chosen flat torsor
chart and put $A=A_\ell/\ell A_\ell$, a nonzero unital
$\F_\ell$-algebra.  Let a Bost pivot have the displayed
expansion.  If $v_\ell(c_n)\ge0$ there is no positive denominator penalty.
Otherwise put
\[
 a=-v_\ell(c_n)>0.
\]
By \cref{lem:pivot-truncation}, $1\le a\le s$.  Put
\[
 \widetilde F_\lambda=f^CF_\lambda,
 \qquad
 \widetilde n=n+C.
\]
For $1\le e\le s$ define the shifted truncated product
\begin{equation}\label{eq:Be-product}
 B_e(f,u,z)=f^CQ_{e,\lambda}(f,u)R_{e,<\ell}(z)
 \in A_\ell[f,z].
\end{equation}
By \cref{lem:pivot-truncation},
\[
 \deg_fB_e\le D+2C<\ell,
 \qquad
 \deg_zB_e<\ell,
 \qquad
 B_e(f,u,0)=0.
\]
Choose coefficientwise representatives
$B_{e,0},\ldots,B_{e,e-1}\in A_\ell[f,z]$ satisfying
\begin{equation}\label{eq:B-digit-expansion}
 B_e\equiv\sum_{\nu=0}^{e-1}\ell^\nu B_{e,\nu}
 \pmod{\ell^eA_\ell[f,z]},
\end{equation}
with the same degree bounds and with $B_{e,\nu}(f,u,0)=0$.  Such digits are
obtained recursively coefficient by coefficient; flatness makes
multiplication by $\ell$ injective in $A_\ell$.

For $1\le b\le s$ define in the special-fibre coefficient algebra
\begin{equation}\label{eq:Gamma-pole-grade}
 \Gamma_b(f,z)=
 \sum_{e\ge b}\overline B_{e,e-b}(f,u,z).
\end{equation}
Let $\Pi_b(f)$ be the pole-grade-$b$ associated-graded symbol of the shifted
truncated expression
\[
 f^CH_\lambda^\flat+
 \sum_{e=1}^s\ell^{-e}B_e(f,u,\varphi_\ell(f)).
\]
Taylor expansion at
$\varphi_\ell(f)=f^\ell+\ell\Delta_\ell(f)$ gives the exact identity
\begin{equation}\label{eq:pole-symbol}
 \boxed{
 \Pi_b(f)=\sum_{t=0}^{s-b}
 \frac{\overline\Delta_\ell(f)^t}{t!}
 \left.\partial_z^t\Gamma_{b+t}(f,z)\right|_{z=f^\ell}.}
\end{equation}
Indeed a digit $B_{e,\nu}$ in the $t$th Taylor term carries the factor
$\ell^{\nu+t-e}$ and hence has pole grade $b=e-\nu-t$; grouping all terms of
fixed $b$ gives \eqref{eq:pole-symbol}.  Since $\ell>s$, every $t!$ occurring
here is a unit.  Notice that digitizing the full product $B_e$, rather than
$Q_{e,\lambda}$ alone, also records positive $\ell$-adic valuations of the
integral coefficients of $R_{e,<\ell}$.

\begin{proposition}[Torsor Cartier strictness]\label{prop:Cartier}
For the actual pole grade $a$, the coefficient of $f^{\widetilde n}$ in
$\Gamma_a(f,f^\ell)$ is nonzero.  Consequently there exist
\[
 -C\le j\le D+C,
 \qquad r\ge1,
\]
such that
\begin{equation}\label{eq:njr}
 n=j+\ell r.
\end{equation}
Hence the positive denominator penalty satisfies
\begin{equation}\label{eq:large-local-height}
 \delta_\ell(\psi_D^{(n)})
 \le s\,v_\ell
 \left([\max\{n-D-C,1\},\ldots,n+C]\right)\log\ell.
\end{equation}
\end{proposition}

\begin{proof}
The shifted scalar series has expansion
\[
 \widetilde F_\lambda
 =c_nf^{\widetilde n}+O(f^{\widetilde n+1}).
\]
For every $k<\widetilde n$ its $f^k$-coefficient is zero.  Taking successive
positive pole-grade symbols of that zero coefficient shows
\[
 \Pi_b(f)\in f^{\widetilde n}A[[f]]
 \qquad(1\le b\le s).
\]
We prove by descending induction on $b$ that
\begin{equation}\label{eq:pole-grade-induction}
 \Gamma_b(f,f^\ell)\in f^{\widetilde n}A[[f]],
 \qquad
 \Pi_b(f)\equiv\Gamma_b(f,f^\ell)
 \pmod{f^{\widetilde n+1}}.
\end{equation}
For $b=s$, \eqref{eq:pole-symbol} has only the term $t=0$, so the assertion
follows.  Suppose it is known for $b+1,\ldots,s$.  For every $t\ge1$, the
induction hypothesis gives
$\Gamma_{b+t}(f,f^\ell)\in f^{\widetilde n}A[[f]]$.  Since
$\deg_f\Gamma_{b+t}<\ell$, \cref{lem:no-backflow} yields
\[
 \frac{\overline\Delta_\ell(f)^t}{t!}
 \left.\partial_z^t\Gamma_{b+t}(f,z)\right|_{z=f^\ell}
 \in f^{\widetilde n+t}A[[f]]
 \subset f^{\widetilde n+1}A[[f]].
\]
The $t=0$ term is $\Gamma_b(f,f^\ell)$, proving the congruence in
\eqref{eq:pole-grade-induction}; because $\Pi_b\in f^{\widetilde n}A[[f]]$,
it also proves the asserted divisibility of $\Gamma_b$.

At the actual pole grade $a$, the coefficient of $f^{\widetilde n}$ in
$\Pi_a$ is the reduction of the unit $\ell^ac_n$.  It is a nonzero scalar in
$\F_\ell$ and remains nonzero in the nonzero unital algebra $A$.  By
\eqref{eq:pole-grade-induction}, the coefficient of $f^{\widetilde n}$ in
$\Gamma_a(f,f^\ell)$ is therefore nonzero.  Some monomial contributing to it
has exponent
\[
 \widetilde n=j'+\ell r,
 \qquad 0\le j'\le D+2C.
\]
Every digit $B_{e,\nu}$ has zero $z$-constant term, so $r\ge1$.  Returning to
the unshifted order and putting $j=j'-C$ gives
\[
 n=j+\ell r,
 \qquad -C\le j\le D+C.
\]
The positive integer $n-j=\ell r$ lies in
$[\max\{n-D-C,1\},n+C]$, so the $\ell$-adic valuation of the lcm of this
interval is at least one.  Since $a\le s$, this proves
\eqref{eq:large-local-height}.
\end{proof}

\begin{remark}[Signed height versus positive penalty]
Bost's local height is the signed quantity
$h^B_\ell=-v_\ell(c_n)\log\ell$.  The local Dwork estimates bound
$\delta_\ell=\max\{0,-v_\ell(c_n)\}\log\ell$, and
$h^B_\ell\le\delta_\ell$.  Replacing the signed height by this upper bound is
therefore legitimate.
\end{remark}

\subsection{The aggregate large-prime sum}

For $m\ge2$ and $1<\xi<m$, define $I^m_\xi(\xi)$ by
\eqref{eq:I-def-intro}.

\begin{lemma}[CDT prime-window summation]\label{lem:CDT-sum}
If $V_D$ satisfies \eqref{eq:zero-estimate}, then
\[
 \sum_{n\in V_D}\sum_{\ell>\xi D}
 v_\ell([\max\{n-D,1\},\ldots,n])\log\ell
 \le I^m_\xi(\xi)D^2+o(D^2).
\]
\end{lemma}

\begin{proof}
This is the one-variable summation in \cite[Section~7.3]{CDT2024}.  A prime
$\ell=xD$ can contribute only if a pivot lies within one multiplier window of
a positive multiple of $\ell$.  The distinct pivots are maximally packed in
an interval of normalized length $m$.  Integrating the resulting packing
function and partitioning at the points $1+k\xi$ gives
\eqref{eq:I-def-intro}.
\end{proof}

\begin{proposition}[Large-prime arithmetic cost]\label{prop:large-cost}
The total residual positive denominator penalty at primes $\ell>N^\sharp$ is
\[
 \mathscr A_D^{\mathrm{large}}
 \le sI^{s+1}_\xi(\xi)D^2+o(D^2).
\]
\end{proposition}

\begin{proof}
Since $N^\sharp=\xi D+O(1)$, the sum over $\ell>N^\sharp$ is bounded
by the CDT sum over $\ell>\xi D$, up to a fixed-width boundary error.
Apply \cref{prop:Cartier,lem:CDT-sum}.  Adding and removing the fixed
$C$-width changes the prime sum by $O_{p,s}(D\log D)$, because only fixed
boundary integers are added for each of $O_s(D)$ pivots.
\end{proof}

\begin{theorem}[Hybrid finite-place theorem]\label{thm:hybrid-local}
For the hybrid source lattice, the finite-place contributions satisfy
\[
 \widehatdeg M_D^{\mathrm{hyb}}
 =\widehatdeg M_D-S_{p,s}(\xi)D^2+o(D^2),
\]
\[
 \mathscr A_D^{\mathrm{small}}=o(D^2),
 \qquad
 \mathscr A_D^{\mathrm{large}}
 \le sI^{s+1}_\xi(\xi)D^2+o(D^2).
\]
\end{theorem}

\begin{proof}
Combine \cref{prop:small-cost,lem:small-residual,prop:large-cost}.
\end{proof}

\section{Bost-slopes assembly}\label{sec:bost}

\begin{proposition}[Bost--CDT leading terms]\label{prop:Bost-leading}
For the unmodified rank-$m$ source,
\begin{equation}\label{eq:baseline-degree}
 \widehatdeg M_D=\frac m2\Energy(\boldsymbol\phi)D^2+o(D^2).
\end{equation}
For the Bost pivots, the nonarithmetic template terms satisfy
\begin{equation}\label{eq:height-leading}
 \sum_{n\in V_D}\sum_vh_v(\psi_D^{(n)})
 \le\left(-\frac{m^2}{2}\LambdaAH(\boldsymbol\phi)
 +m\Energy(\boldsymbol\phi)\right)D^2
 +\mathscr A_D+o(D^2).
\end{equation}
\end{proposition}

\begin{proof}
For the split scalar source, the scalar source-degree asymptotic is
\cite[(7.3.7)]{CDT2024}.  In the adelic normalization,
\cite[Section~2.3 and the proof of Theorem~2.6]{CDT2025} derives the same
leading term from the arithmetic Hilbert--Samuel theorem together with the
Bost--Charles overflow identities
\cite[Theorem~5.4.1 and Proposition~5.4.2]{BostCharles2022}.
The fixed-bundle comparison \cref{lem:fixed-bundle} therefore gives
\eqref{eq:baseline-degree}.

For a pivot of order $n$, the constituent CDT estimates have nonarithmetic
part
\[
 -n\LambdaAH(\boldsymbol\phi)+D\Energy(\boldsymbol\phi),
\]
while the finite coefficient-arithmetic term is the quantity denoted
$\mathscr A_D$ here; see \cite[(7.3.14)--(7.3.15)]{CDT2024}.  There are
$mD+O(1)$ distinct nonnegative pivots and, by
\cref{cor:zero-estimate}, each is $O(D)$.  Hence
\[
 \sum_{n\in V_D}n\ge\binom{\#V_D}{2}
 =\frac12m^2D^2+o(D^2),
\]
and the uniform lower-order terms sum to $o(D^2)$.  Summing the CDT estimates
and retaining $\mathscr A_D$ separately gives \eqref{eq:height-leading}; this
is the leading-coefficient passage in
\cite[(7.3.14)--(7.3.16)]{CDT2024}, with the adelic normalization above.
\end{proof}

\begin{proof}[Proof of \cref{thm:hybrid-holonomy}]
Bost's slopes inequality, in the form of
\cite[(7.3.16)]{CDT2024}, gives
\[
 \widehatdeg M_D^{\mathrm{hyb}}
 \le\sum_{n\in V_D}\sum_vh_v(\psi_D^{(n)}).
\]
By \cref{thm:hybrid-local,prop:Bost-leading},
\[
 \frac m2\Energy-S_{p,s}(\xi)
 \le-\frac{m^2}{2}\LambdaAH+m\Energy
 +sI^{s+1}_\xi(\xi).
\]
Rearranging and using $m=s+1$ gives
\[
 m\left(\LambdaAH-
 \frac{2}{m^2}(S_{p,s}(\xi)+sI^{s+1}_\xi(\xi))\right)
 \le\Energy,
\]
which is \eqref{eq:tau-main}.
\end{proof}

\section{Uniform analytic continuation}\label{sec:continuation}

Put
\[
 a_p=\frac{12}{p-1},
 \qquad
 R_p=p^{a_p}.
\]

\begin{theorem}[Uniform continuation]\label{thm:continuation}
For $p\in\{2,3,5,7\}$ and odd $s\ge3$, the distinguished form
$K_{p,s,\zeta_p(s)/2}$ extends as a rigid analytic section of
$\omega^{1-s}$ to
\[
 U_p=\{\abs{f_p}_p<R_p\}.
\]
The full Eisenstein--Eichler horizontal leaf extends meromorphically over the
same wide-open.  At every finite place $\ell\ne p$, under
$\zeta_p(s)\in\Q$, it is meromorphic on $\abs{f_p}_\ell<1$.  For every
$Y>0$, its complex pullback under $z\mapsto f_p(e^{-2\pi Y}z)$ is
meromorphic near the closed unit disc.
\end{theorem}

\begin{proof}
Input (E1) makes $K$ overconvergent and \cref{lem:Up} gives $U_p$-eigenvalue
one.  The weight $1-s$ is an integer, and Buzzard's theorem imposes no
nonnegativity condition on that integer weight.  Pull to one fixed neat tame
level, apply \cite[Theorem~5.2(2)]{Buzzard2003}, and descend by uniqueness of
analytic continuation.  Thus $K$ extends to the wide-open $W_0(p)$ containing the
cusp $\infty$.

For $p=2,3$, Calegari identifies the two ordinary components and the
intervening supersingular region in the coordinate $f_p$, giving
$W_0(p)=\{\abs{f_p}_p<R_p\}$; see \cite[Section~3]{Calegari2005}.  We spell
out the corresponding calculation for $p=5,7$.  Put $t_p=f_p^{-1}$.
For $p=5$, Kilford's equation \cite[(3)]{Kilford2008} is
\[
 j=\frac{(t_5^2+250t_5+3125)^3}{t_5^5}.
\]
Writing $v=v_5(t_5)$, its Newton-polygon valuation gives
\[
 v_5(j)=v\quad(0<v<5/2),
 \qquad
 v_5(j)=15-5v\quad(5/2<v<3).
\]
Since $j=0$ is the unique supersingular $j$-invariant in characteristic $5$,
the supersingular annulus is $0<v_5(t_5)<3$.  The cusp $\infty$ lies on the
side where $t_5$ has a pole; hence
\[
 W_0(5)=\{v_5(t_5)<3\}=\{\abs{f_5}_5<5^3\}.
\]
For $p=7$, Kilford--McMurdy give $w_7^*t_7=49/t_7$ and, from their explicit
$j$-models, identify the unique supersingular annulus as
$0<v_7(t_7)<2$; see \cite[Equations~(1)--(4) and Section~2]{KilfordMcMurdy2012}.
Thus
\[
 W_0(7)=\{v_7(t_7)<2\}=\{\abs{f_7}_7<7^2\}.
\]

The Gauss--Manin connection is algebraic on the de Rham frame torsor, so
applying it finitely many times does not shrink the base domain.  This extends
the BGG leaf.  At $\ell\ne p$, the raw coefficients have at worst logarithmic
$\ell$-adic denominator growth and hence radius at least one; the integral
inverse coordinate transfers this to the $f$-disc.  Over $\C$, the divisor-sum
coefficients have subexponential growth, giving the final claim.
\end{proof}

\begin{remark}[Strict exhaustion]
The $p$-adic continuation domain is open.  In the holonomy argument one uses
closed affinoids of radius $p^{a_p-\varepsilon}$ and lets
$\varepsilon\downarrow0$.  No boundary convergence is assumed.
\end{remark}

\section{Exact modular collision energy}\label{sec:collision}

For $0<Y<1/p$, define
\[
 \phi_{p,Y}(z)=f_p(e^{-2\pi Y}z)
\]
and $C_p(Y)$ by \eqref{eq:C-intro}.

\begin{proposition}[Modular Jensen formula]\label{prop:Jensen}
One has
\begin{equation}\label{eq:Jensen}
 \iint_{\T^2}
 \log\abs{\phi_{p,Y}(z)-\phi_{p,Y}(w)}\,\dd m(z)\dd m(w)
 =-2\pi Y+C_p(Y).
\end{equation}
\end{proposition}

\begin{proof}
Fix $\tau_x=x+iY$.  Since $f_p$ is a Hauptmodul, the nonidentity zeros of
$f_p(q')-f_p(e^{2\pi i\tau_x})$ are parametrized by cosets
$\Gamma_\infty\backslash\Gamma_0(p)$, equivalently primitive bottom rows
$(c,d)$ with $c>0$ and $p\mid c$.  For
$\gamma=\left(\begin{smallmatrix}a&b\\c&d\end{smallmatrix}\right)$,
\[
 \Im(\gamma\tau_x)=
 \frac{Y}{(cx+d)^2+c^2Y^2}.
\]
Jensen's formula contributes
$2\pi(\Im(\gamma\tau_x)-Y)_+$.  For a fixed $c$, unfold the sum over
primitive $d$:
\[
 \sum_{\substack{d\in\Z\\(c,d)=1}}
 \int_0^1F(cx+d)\,\dd x
 =\frac{\varphi(c)}{c}\int_\R F(u)\,\dd u.
\]
The positive part is nonempty exactly when $cY<1$.  Evaluating the elementary
integral gives
\[
 \frac{4\pi\varphi(c)}{c^2}
 \left(\arccos(cY)-cY\sqrt{1-c^2Y^2}\right).
\]
Only finitely many $c$ occur.  The eta-product factors in $f_p$ have
logarithmic mean zero, while the leading $q$ contributes $-2\pi Y$.
Elliptic stabilizers occur only on a measure-zero set and are accounted for by
the standard orbifold multiplicity.
\end{proof}

At the distinguished $p$-adic place, strict exhaustion of the radius $R_p$
gives a scaling width $12\log p/(p-1)$.  Combining it with
\cref{prop:Jensen} yields
\begin{equation}\label{eq:Lambda-energy}
 \Lambda_p(Y)=\frac{12\log p}{p-1}-2\pi Y,
 \qquad
 \Energy_p(Y)=\Lambda_p(Y)+C_p(Y).
\end{equation}

\section{The twenty-two applications}\label{sec:applications}

For each target we use exact rational witnesses $\xi,Y$ listed in
\cref{tab:22}.  Every witness satisfies
\[
 1<\xi<s+1,
 \qquad c_p\xi<1,
 \qquad0<Y<1/p.
\]
The value of $I^{s+1}_\xi(\xi)$ and hence $\tau_{p,s}(\xi)$ is rational.
The collision sum has exactly the finitely many active layers
$c=p,2p,\ldots$ with $cY<1$.

\begingroup\small
\begin{longtable}{@{}rrrrrl@{}}
\caption{Exact rational witnesses and independently certified lower margins.}
\label{tab:22}\\
\toprule
$p$&$s$&$\xi$&$\tau_{p,s}(\xi)$&$Y$& certified lower bound\\
\midrule
\endfirsthead
\toprule
$p$&$s$&$\xi$&$\tau_{p,s}(\xi)$&$Y$& certified lower bound\\
\midrule
\endhead
2 & 3 & $22/19$ & $1321/608$ & $16/79$ & $9.745786633966$ \\
2 & 5 & $29/26$ & $18445/5616$ & $6/49$ & $13.454727090366$ \\
2 & 7 & $110/101$ & $105585/25856$ & $1/11$ & $15.805524331097$ \\
2 & 9 & $89/83$ & $10906993/2324000$ & $3/43$ & $17.190005995811$ \\
2 & 11 & $262/247$ & $464385991/89631360$ & $3/52$ & $17.790824360622$ \\
2 & 13 & $181/172$ & $144989487/25958240$ & $2/41$ & $17.737379764768$ \\
2 & 15 & $478/457$ & $16666007411/2810615808$ & $3/71$ & $17.124944575766$ \\
2 & 17 & $305/293$ & $12175769323/1954839744$ & $3/80$ & $16.023242852853$ \\
2 & 19 & $758/731$ & $11693630797/1801184000$ & $1/30$ & $14.484310655378$ \\
2 & 21 & $461/446$ & $536638388131/79764338368$ & $1/33$ & $12.553794492186$ \\
2 & 23 & $1102/1069$ & $22108077629089/3185347156992$ & $1/36$ & $10.264641765760$ \\
2 & 25 & $649/631$ & $21299386717285/2985462031552$ & $1/39$ & $7.646580292916$ \\
2 & 27 & $1510/1471$ & $456079886220493/62372115303680$ & $1/42$ & $4.724811329943$ \\
2 & 29 & $869/848$ & $14321093408042087/1915475381820000$ & $1/46$ & $1.522247623229$ \\
3 & 3 & $33/29$ & $2029/928$ & $4/27$ & $6.200294149415$ \\
3 & 5 & $87/79$ & $112477/34128$ & $7/73$ & $6.941905556192$ \\
3 & 7 & $55/51$ & $53435/13056$ & $1/15$ & $6.163732078651$ \\
3 & 9 & $267/251$ & $33037241/7028000$ & $1/19$ & $4.332047506141$ \\
3 & 11 & $393/373$ & $702154669/135354240$ & $3/70$ & $1.679644254174$ \\
5 & 3 & $11/10$ & $177/80$ & $7/71$ & $2.515479614080$ \\
5 & 5 & $29/27$ & $12899/3888$ & $1/16$ & $0.131799356827$ \\
7 & 3 & $33/31$ & $2219/992$ & $1/14$ & $0.436413057785$ \\
\bottomrule
\end{longtable}
\endgroup

\begin{proof}[Proof of \cref{thm:main}]
Fix one row of \cref{tab:22} and suppose $\zeta_p(s)\in\Q$.  By
\cref{thm:continuation,prop:Jensen}, the templates have width and energy
\eqref{eq:Lambda-energy}.  Apply \cref{thm:hybrid-holonomy} with the displayed
$\xi$.  It forces
\[
 (s+1)(\Lambda_p(Y)-\tau_{p,s}(\xi))
 \le\Lambda_p(Y)+C_p(Y),
\]
or equivalently
\[
 \mathfrak M_{p,s}(\xi,Y)
 =s\Lambda_p(Y)-(s+1)\tau_{p,s}(\xi)-C_p(Y)\le0.
\]
The interval certificate gives a strictly positive lower bound in every row,
a contradiction.
\end{proof}

\section{Logical scope and dependency audit}\label{sec:audit}

The proof is intentionally organized so that every delicate interface can be
checked independently.

\begingroup\small
\begin{longtable}{@{}>{\raggedright\arraybackslash}p{0.23\textwidth}>{\raggedright\arraybackslash}p{0.35\textwidth}>{\raggedright\arraybackslash}p{0.34\textwidth}@{}}
\toprule
Gate & Source & Use in the proof\\
\midrule
Calegari Eisenstein family & \cite[Theorem~2.3 and Lemma~2.4]{Calegari2005}
& Identifies $K_{p,s}$ and gives overconvergence.\\
Buzzard continuation & \cite[Theorem~5.2]{Buzzard2003}
& Extends the finite-slope form to $W_0(p)$.\\
Hasse invariant & \cite{Katz1973} plus classical Eisenstein congruences
& Algebraizes $x_\sigma$ modulo an auxiliary prime and bounds its degree.\\
De Rham frame torsor & \cite[Propositions~A.14--A.15]{GPRJ2025}
& Supplies the neat-level torsor, filtration, weight decomposition, and
Gauss--Manin structure.  Descent and spreading out are internal.\\
Bost/CDT & \cite{Bost2001};
\cite[Theorem~5.4.1 and Proposition~5.4.2]{BostCharles2022};
\cite[Theorem~3.2.13 and (7.3.7), (7.3.14)--(7.3.16)]{CDT2024};
\cite[Theorem~2.6]{CDT2025}
& Supplies the slopes inequality, overflow/source identities, scalar zero
estimate, pivot-height bounds, adelic normalization, and prime-window
summation; \cref{lem:fixed-bundle} proves the fixed-bundle norm comparison.\\
BGG extension & \cite[Proposition~2.4.3]{RodriguezCamargo2023};
\cite[Proposition~2.4.1 and Remark~2.4.2]{Urban2013}; internal
\cref{lem:BGG,prop:split-extension}
& Fixes the exact global BGG--Bol normalization and constructs the algebraic
Eisenstein--Eichler connection.\\
Small-prime arithmetic & Internal,
\cref{prop:single-unipotent,thm:genuine-Hasse,lem:actual-deepest-row,prop:small-cost}
& Proves the normalized deepest-row degree bound, genuine-source Hasse
rank--nullity saving, fixed-width cutoff control, and determinant cost.\\
Large-prime arithmetic & Internal,
\cref{lem:no-backflow,lem:pivot-truncation,prop:Cartier}
& Truncates before taking pole grades, digitizes each full truncated product,
and gives the torsor-valued Cartier window without scalarization.\\
Collision energy & Internal, \cref{prop:Jensen}
& Gives the exact complex Bost--Charles energy.\\
Numerics & Appendix~\ref{app:certificate}
& Gives exact rational witnesses and outward interval enclosures.\\
\bottomrule
\end{longtable}
\endgroup

The attribution boundary is as follows.  GPRJ supplies the torsor and its
Gauss--Manin structures at neat level; \cref{lem:neat-descent} proves the
descent and integral spreading used here.  CDT supplies scalar source and
pivot estimates; \cref{lem:fixed-bundle} proves that passage to the fixed
vector-bundle source changes only $o(D^2)$.  Katz supplies the Hasse invariant
and $q$-expansion principle; \cref{prop:single-unipotent} proves the
single-action degree bound for the genuine deepest multiplier, and the
codimension estimate is the internal rank--nullity argument of
\cref{thm:genuine-Hasse}.

\begin{warning}[Obsolete proof routes]
The manuscript does not inherit the global scalar-gauge assertion from older
prime-window drafts.  A nonzero-weight Tate trivialization transforms by a
nontrivial power of $c\tau+d$ and therefore cannot equal a rational function
of the Hauptmodul.  The Rees--Bol matrix controls the unipotent coordinate,
not the diagonal torus.  This is why the large-prime proof stays on the frame
torsor and the small-prime proof uses only a mod-$\ell$ Hasse comparison.
\end{warning}

\appendix

\section{Independent rational fixed-point interval certificate}
\label{app:certificate}

The supplementary program
\path{hybrid_rational_interval_certificate.py} reconstructs every numerical
quantity independently from the formulas of the paper.  It uses only Python's
standard-library arbitrary-precision integers and \texttt{Fraction}; no binary
floating point or external interval package enters the certification.
All algebraic quantities are exact.  Real intervals are represented by integer
endpoints at the fixed scale $10^{90}$, and every arithmetic operation is
rounded outward.

The transcendental constants are enclosed from elementary convergent series
with explicit remainder bounds:
\[
 \pi=16\arctan\frac15-4\arctan\frac1{239},
 \qquad
 \log p=2\operatorname{arctanh}\frac{p-1}{p+1},
\]
\[
 \arctan z=\sum_{k\ge0}\frac{(-1)^kz^{2k+1}}{2k+1},
 \qquad
 \arccos x=2\arctan\sqrt{\frac{1-x}{1+x}}.
\]
The alternating arctangent remainder is bounded by the first omitted term;
the positive arctanh tail is bounded by a geometric series.  Square roots of
rationals are enclosed by exact integer square root at scale $10^{90}$.
Thus every printed interval follows from integer inequalities that can be
checked independently of a transcendental-function library.

For each pair $(p,s)$ the script computes
\[
 \xi=\frac{12(s(s+1)-1)}{12s^2+(s-1)(p+1)},
\]
then evaluates $K_m(\xi)$, $I^m_\xi(\xi)$, $S_{p,s}(\xi)$, and
$\tau_{p,s}(\xi)$ exactly.  The active collision layers are selected by the
exact rational test $cY<1$.  The program asserts that the lower endpoint of
every one of the twenty-two final intervals is positive.  Its smallest case is
\[
\begin{aligned}
 \mathfrak M_{5,5}\left(\frac{29}{27},\frac1{16}\right)
 \in [&0.131799356827016832557664457131479890735826671587851068262655,\\
      &0.131799356827016832557664457131479890735826671587851068262656].
\end{aligned}
\]

For reproducibility, the accompanying manifest \path{hybrid_rational_interval_certificate_sha256.txt} records the SHA-256 hashes
of the certificate script and its output:
\begin{verbatim}
1408c9092d0c34917253c8f520db56853112c58640d15c98d58b76a61a0478f5
  hybrid_rational_interval_certificate.py
9e1d8f74eb198c7e7b0c7420a1e8021d7bad4d2bc8014cc939b6b353111c0da0
  hybrid_rational_interval_certificate_output.txt
\end{verbatim}

\section{Explicit prime-window function}\label{app:window}

For completeness, let $m\ge2$, $1<\xi<m$, and
$K=\floor{(m-1)/\xi}$.  The normalized pivot interval has length $m$ and a
large prime $\ell=xD$ contributes when a pivot lies in a degree-$D$ interval
ending at a positive multiple of $\ell$.  Maximally packing distinct pivots
and integrating the resulting triangular boundary term gives
\[
 I^m_\xi(\xi)=
 \left(m-\frac12\right)H_K-K\xi
 +\frac{(m-(K+1)\xi)_+^2}{2(K+1)}.
\]
This is the specialization of the CDT summation used in
\cref{lem:CDT-sum}.

\begin{thebibliography}{99}

\bibitem{Bost2001}
J.-B.~Bost,
\newblock Algebraic leaves of algebraic foliations over number fields,
\newblock \emph{Publ. Math. Inst. Hautes \`Etudes Sci.} \textbf{93} (2001),
161--221.

\bibitem{BostCharles2022}
J.-B.~Bost and F.~Charles,
\newblock \emph{Quasi-projective and formal-analytic arithmetic surfaces},
\newblock Annals of Mathematics Studies \textbf{224}, Princeton University
Press, Princeton, NJ, 2026; arXiv:2206.14242.

\bibitem{Buzzard2003}
K.~Buzzard,
\newblock Analytic continuation of overconvergent eigenforms,
\newblock \emph{J. Amer. Math. Soc.} \textbf{16} (2003), 29--55;
\newblock doi:10.1090/S0894-0347-02-00405-8.

\bibitem{Calegari2005}
F.~Calegari,
\newblock Irrationality of certain $p$-adic periods for small $p$,
\newblock \emph{Int. Math. Res. Not.} 2005, no.~20, 1235--1249;
\newblock doi:10.1155/IMRN.2005.1235.

\bibitem{CDT2024}
F.~Calegari, V.~Dimitrov and Y.~Tang,
\newblock The linear independence of $1$, $\zeta(2)$, and
$L(2,\chi_{-3})$,
\newblock arXiv:2408.15403v2 (2024).

\bibitem{CDT2025}
F.~Calegari, V.~Dimitrov and Y.~Tang,
\newblock Arithmetic holonomy bounds and effective Diophantine approximation,
\newblock in \emph{Proceedings of the International Congress of
Mathematicians 2026}, vol.~3, SIAM, 2026, pp.~356--376;
\newblock doi:10.1137/25M1806776; arXiv:2510.04156.

\bibitem{GPRJ2025}
A.~Graham, V.~Pilloni and J.~Rodrigues Jacinto,
\newblock $p$-adic interpolation of Gauss--Manin connections on nearly
overconvergent modular forms and $p$-adic $L$-functions,
\newblock \emph{Compos. Math.} \textbf{161} (2025), 2380--2441;
\newblock doi:10.1112/S0010437X25102479.

\bibitem{Katz1973}
N.~M.~Katz,
\newblock $p$-adic properties of modular schemes and modular forms,
\newblock in \emph{Modular Functions of One Variable III},
Lecture Notes in Math. \textbf{350}, Springer, 1973, pp.~69--190.

\bibitem{Kilford2008}
L.~J.~P.~Kilford,
\newblock On the slopes of the $U_5$ operator acting on overconvergent modular forms,
\newblock \emph{J. Th\'eor. Nombres Bordeaux} \textbf{20} (2008), no.~1,
165--182; doi:10.5802/jtnb.620.

\bibitem{KilfordMcMurdy2012}
L.~J.~P.~Kilford and K.~McMurdy,
\newblock Slopes of the $U_7$ operator acting on a space of overconvergent modular forms,
\newblock \emph{LMS J. Comput. Math.} \textbf{15} (2012), 113--139;
\newblock doi:10.1112/S1461157012000095.

\bibitem{Lai2025}
L.~Lai,
\newblock On the irrationality of certain $2$-adic zeta values,
\newblock \emph{Int. J. Number Theory} \textbf{21} (2025), no.~1, 207--235;
\newblock doi:10.1142/S1793042125500113.

\bibitem{LaiLupuSprang2025}
L.~Lai, C.~Lupu and J.~Sprang,
\newblock On the irrationality of certain $p$-adic zeta values,
\newblock \emph{Res. Math. Sci.} \textbf{12} (2025), article~77;
\newblock doi:10.1007/s40687-025-00559-x; arXiv:2505.23088.

\bibitem{LaiSprang2025}
L.~Lai and J.~Sprang,
\newblock Many $p$-adic odd zeta values are irrational,
\newblock \emph{Michigan Math. J.}, advance publication (2026), 1--36;
\newblock doi:10.1307/mmj/20236490; arXiv:2306.10393v2.

\bibitem{RodriguezCamargo2023}
J.~E.~Rodr\'iguez Camargo,
\newblock $p$-adic Eichler--Shimura maps for the modular curve,
\newblock \emph{Compos. Math.} \textbf{159} (2023), no.~6, 1214--1249;
\newblock doi:10.1112/S0010437X23007182; arXiv:2102.13099.

\bibitem{Urban2013}
E.~Urban,
\newblock Nearly overconvergent modular forms,
\newblock in \emph{Iwasawa Theory 2012}, Contributions in Mathematical and
Computational Sciences \textbf{7}, Springer, 2014, pp.~401--441;
\newblock doi:10.1007/978-3-642-55245-8\_14.

\end{thebibliography}

\end{document}
